% Created 2020-11-21 Sat 13:36
% Intended LaTeX compiler: pdflatex
\documentclass[11pt]{article}
\usepackage[utf8]{inputenc}
\usepackage[T1]{fontenc}
\usepackage{graphicx}
\usepackage{grffile}
\usepackage{longtable}
\usepackage{wrapfig}
\usepackage{rotating}
\usepackage[normalem]{ulem}
\usepackage{amsmath}
\usepackage{textcomp}
\usepackage{amssymb}
\usepackage{capt-of}
\usepackage{hyperref}
\author{Pascal Fleury}
\date{\today}
\title{Paf's portable Emacs configuration}
\hypersetup{
 pdfauthor={Pascal Fleury},
 pdftitle={Paf's portable Emacs configuration},
 pdfkeywords={},
 pdfsubject={},
 pdfcreator={Emacs 26.3 (Org mode 9.4)},
 pdflang={English}}
\begin{document}

\maketitle
\tableofcontents

This is my Emacs config. You need to tangle this into a file that then
gets loaded by Emacs: \texttt{C-c C-v t} [org-babel-tangle]. Below, I also
explain how this tangling is automated.

Find extensive documentation about how to do this \href{https://github.com/larstvei/dot-emacs}{here}.

\section{Personal plans}
\label{sec:org613b8bc}
I regularly try out new packages, this is my current list of things
being evaluated.  They usually live temporarily in my \texttt{\textasciitilde{}/.emacs} until
I am happy, in which case I move their config into this file so that
it gets replicated on all machines I work on with Emacs.

\subsection{Currently under review}
\label{sec:org2cd2f37}
\begin{itemize}
\item \href{https://github.com/Kungsgeten/yankpad}{yankpad} + \href{https://github.com/joaotavora/yasnippet}{yasnippet}
\item \href{https://github.com/yjwen/org-reveal}{ox-reveal} (see \href{https://github.com/yjwen/org-reveal\#set-the-location-of-revealjs}{installation})
\end{itemize}

\section{One-time initial setup}
\label{sec:org9901551}
I have my config in directory \texttt{\textasciitilde{}/Emacs} which is where I clone this repository. The config setup is maintained purely in the \texttt{\textasciitilde{}/Emacs/emacs\_setup.org} file.

In your \texttt{\textasciitilde{}/.emacs} file, all you need to add is

\begin{verbatim}
;; Loads PAF's emacs setup with bootstrap
(load-file "~/Emacs/emacs_setup.el")
\end{verbatim}

\subsection{Bootstrap}
\label{sec:org6a07b41}
Initially when cloning this repository, you have the \texttt{emacs\_setup.org}
file, that contains the config that you adapt to your specific setup,
and an \texttt{emacs\_setup.el} with a bootstrap content that will tangle and
compile the org file, \emph{and replace itself}. This is useful the very
first time.

After that, the config itself should have the hook to re-tangle and
re-compile the setup at each save.

Therefore my setup is very easy to install, and it needs these steps:

\begin{enumerate}
\item clone this repo into \texttt{\textasciitilde{}/Emacs}
\item add the one line in you \texttt{\textasciitilde{}/.emacs}
\item make sure Emacs re-interprets its init (you could restart it)
\end{enumerate}

It may be that \href{https://github.com/jwiegley/use-package}{use-package} is not installed on your setup, so it will
first try to install that. After that step, it will also start
installing any package that is marked as needed in this config
automatically.

The original content of the \texttt{emacs\_setup.el} is as follows:

\begin{verbatim}
;; This is the initial state of the file to be loaded.
;; It will replace itself with the actual configuration at first run.

(require 'org) ; We can't tangle without org!

(setq config_base (expand-file-name "emacs_setup"
                                    (file-name-directory
                                     (or load-file-name buffer-file-name))))
(find-file (concat config_base ".org"))        ; Open the configuration
(org-babel-tangle)                             ; tangle it
(load-file (concat config_base ".el"))         ; load it
(byte-compile-file (concat config_base ".el")) ; finally byte-compile it
\end{verbatim}

\subsection{Recompile all packages}
\label{sec:orgac09f84}
This will force-recompile everything in \texttt{\textasciitilde{}/.emacs.d/elpa/...} Just run \texttt{M-:} and then enter this:
\begin{verbatim}
(byte-recompile-directory package-user-dir nil 'force)
\end{verbatim}
or simply \texttt{C-x C-e} at the end of that line.

\subsection{One-time configure}
\label{sec:org7e4fe16}
To preserve the original state of this file when updating the git repos with new config settings, execute the following block once (\texttt{C-c C-c}):

\begin{verbatim}
#!/bin/bash
# Make git ignore the tangled & updated emacs_setup.el
if [[ -z "$(which git)" ]]; then
  echo "You will need 'git' to be installed !"
  exit 1
fi
if  [[ -z "$(which emacs)" ]]; then
  echo "You might need 'emacs' for this to be useful !"
  exit 1
fi

GIT_ROOT=$(dirname $0)
(cd ${GIT_ROOT} && git update-index --skip-worktree emacs_setup.el)

# Maybe this is a new install, .emacs does not exist
test -e ~/.emacs || touch ~/.emacs

# Initial tangle of files.
emacs --batch --load ${GIT_ROOT}/emacs_setup.el

# Add the load-file as the first thing in the user's ~/.emacs
# If not yet added.
declare lines=$(grep ';; dot_emacs.el' ~/.emacs | wc -l)
if (( lines < 1 )); then
  echo "Added loading the config in your ~/.emacs"
  echo ";; dot_emacs.el" > ~/.emacs.new
  cat ${GIT_ROOT}/dot_emacs.el >> ~/.emacs.new
  cat ~/.emacs >> ~/.emacs.new
  mv ~/.emacs.new ~/.emacs
else
  echo "Config in your ~/.emacs already set up!"
fi

# Install system dependencies
echo "Installing dependencies"
bash ${GIT_ROOT}/install_deps.sh

echo "Cleanup"
rm ${GIT_ROOT}/dot_emacs.el
rm ${GIT_ROOT}/install_deps.sh
\end{verbatim}

This script is then used to install the needed packages on the system.
\begin{verbatim}
#!/bin/bash

# Trick to make it work on Termux
which "ls" || pkg install debianutils

# This is a bit of heuristics to find out what the install system is
# They are attempted in this order, put the least likely first.
declare -a PKG_MGRS=("pkg" "brew" "apt-get")

PKG_PREFIX_apt_get="sudo"

for pkg in "${PKG_MGRS[@]}"; do
    if [[ -x "$(which ${pkg})" ]]; then
        INSTALLER="${pkg}"
        break
    fi
done
if [[ -z "${INSTALLER}" ]]; then
    echo "Did not find a suitable installer (tried ${PKG_MGRS[@]})"
    exit 1
fi

# This is the function to call to install anything. It can optionally
# check for a binary and avoid installing if it's found.  install_pkg
# [-x <binary>] <package>
function install_pkg() {
    if [[ "$1" == "-x" ]]; then
        local binary="$(which $2)"
        if [[ -n "${binary}" && -x "${binary}" ]]; then
            echo "Found $2 (${binary}), nothing to install for $3."
            return
        fi
        shift 2
    fi

    local token=$(echo -n ${INSTALLER} | tr -c '0-9a-zA-Z_' '_')
    local prefix_var="PKG_PREFIX_${token}"

    echo "Trying: ${INSTALLER} install $*"
    ${!prefix_var} $(which ${INSTALLER}) install "$@"
}
\end{verbatim}

\section{Initialize Emacs}
\label{sec:org4b8b45a}
This section sets up Emacs so it can tangle the config, find \texttt{use-package}, and find the ELPA repositories where to get the new packes from.
\subsection{Info header}
\label{sec:org33d6f8e}
Just to add a little information in the tangled file.
\begin{verbatim}
;; ===== this file was auto-tangled, only edit the emacs_setup.org =====
\end{verbatim}

\subsection{melpa}
\label{sec:org06a58ef}
Make sure we have the package system initialized before we load anything.
\begin{verbatim}
(require 'package)
(when (< emacs-major-version 27)
  (package-initialize))
\end{verbatim}

Adding my choice of packages repositories.
\#+NAME melpa-setup
\begin{verbatim}
(setq package-archives '(("org" . "https://orgmode.org/elpa/")
                         ("stable-melpa" . "https://stable.melpa.org/packages/")
                         ("melpa" . "https://melpa.org/packages/")
                         ("gnu" . "https://elpa.gnu.org/packages/")
                         ; ("marmalade" . "https://marmalade-repo.org/packages/")
                        ))
\end{verbatim}
\subsection{use-package}
\label{sec:org80b4982}
I use \texttt{use-package} for most configuration, and that needs to be at the top of the file.  \texttt{use-package} verifies the presence of the requested package, otherwise installs it, and presents convenient sections for configs of variables, key bindings etc. that happen only if the package is actually loaded.

First, make sure it gets installed if it is not there yet.
\begin{verbatim}
;; make sure use-package is installed
(unless (package-installed-p 'use-package)
  (package-refresh-contents)
  (package-install 'use-package))
(require 'use-package)
\end{verbatim}

\begin{verbatim}
(eval-when-compile (require 'use-package))
\end{verbatim}
\subsection{tangle-this-config}
\label{sec:org6715b1e}
I set this up to tangle the init org-mode file into the actual Emacs init file as soon as I save it.
\begin{verbatim}
(defun tangle-init ()
  "If the current buffer is 'init.org' the code-blocks are
tangled, and the tangled file is compiled."
  (when (equal (buffer-file-name)
               (expand-file-name "~/Emacs/emacs_setup.org"))
    ;; Avoid running hooks when tangling.
    (let ((prog-mode-hook nil))
      (org-babel-tangle)
      (byte-compile-file "~/Emacs/emacs_setup.el"))))

(add-hook 'after-save-hook 'tangle-init)
\end{verbatim}
\section{Personal initialization}
\label{sec:org26342f2}
\subsection{Clear \texttt{C-c f} so I can use it as a prefix}
\label{sec:org3b3a635}
Remove \texttt{C-c f} that I want to use for \textbf{me} personally as a prefix.
\begin{verbatim}
(global-set-key (kbd "C-c f") nil)
\end{verbatim}

\section{Helper functions}
\label{sec:org592ef9d}
\subsection{add-hook-run-once}
\label{sec:org0127b62}
Use instead of add-hook to run it a single time.
\href{https://emacs.stackexchange.com/questions/3323/is-there-any-way-to-run-a-hook-function-only-once}{found here}
\begin{verbatim}
(defmacro add-hook-run-once (hook function &optional append local)
  "Like add-hook, but remove the hook after it is called"
  (let ((sym (make-symbol "#once")))
    `(progn
       (defun ,sym ()
         (remove-hook ,hook ',sym ,local)
         (funcall ,function))
       (add-hook ,hook ',sym ,append ,local))))
\end{verbatim}
\section{Environment}
\label{sec:org45c7b6e}
\subsection{Browser default}
\label{sec:org0845682}
\begin{verbatim}
(setq browse-url-generic-program (executable-find "google-chrome")
  browse-url-browser-function 'browse-url-generic)
\end{verbatim}
\subsection{Setup server}
\label{sec:org2ece2bb}
Start the background server, so we can use emacsclient.
\begin{verbatim}
(server-start)
\end{verbatim}
\subsection{UTF-8}
\label{sec:orgcbb2205}
Make Emacs request UTF-8 first when pasting stuff.
\begin{verbatim}
(use-package unicode-escape
  :ensure t
  :init
  (setq x-select-request-type '(UTF8_STRING COMPOUND_TEXT TEXT STRING)))
(set-language-environment "UTF-8")
\end{verbatim}
\subsection{Newline (only Unix wanted)}
\label{sec:org7d0a122}
This should automatically convert any files with dos or Mac line endings into Unix style ones. Code found \href{https://www.emacswiki.org/emacs/EndOfLineTips}{here}.
\begin{verbatim}
(defun no-junk-please-we-are-unixish ()
  (let ((coding-str (symbol-name buffer-file-coding-system)))
    (when (string-match "-\\(?:dos\\|mac\\)$" coding-str)
      (set-buffer-file-coding-system 'unix))))

(add-hook 'find-file-hook 'no-junk-please-we-are-unixish)
\end{verbatim}

\subsection{save \& restore buffers}
\label{sec:orge43c89b}

Turn on global:
\begin{verbatim}

(setq desktop-path (list "~/.emacs.d/savehist"))
(setq desktop-dirname "~/.emacs.d/savehist")
(setq desktop-restore-eager 5)
(setq desktop-load-locked-desktop t)
(desktop-save-mode 1)

(savehist-mode 1)
(setq history-length t)
(setq history-delete-duplicates t)
(setq savehist-save-minibuffer-history 1)
(setq savehist-additional-variables '(kill-ring search-ring regexp-search-ring))

\end{verbatim}

But skip the followings:
\begin{verbatim}
(setq desktop-buffers-not-to-save
     (concat "\\("
             "^nn\\.a[0-9]+\\|\\.log\\|(ftp)\\|^tags\\|^TAGS"
             "\\|\\.emacs.*\\|\\.diary\\|\\.newsrc-dribble\\|\\.bbdb"
             "\\)$"))
(add-to-list 'desktop-modes-not-to-save 'dired-mode)
(add-to-list 'desktop-modes-not-to-save 'Info-mode)
(add-to-list 'desktop-modes-not-to-save 'info-lookup-mode)
(add-to-list 'desktop-modes-not-to-save 'fundamental-mode)
\end{verbatim}

\subsection{auto revert}
\label{sec:org393c6e1}
Use \texttt{auto-revert}, which reloads a file if it's updated on disk
and not modified in the buffer.
\begin{verbatim}
(global-auto-revert-mode 1)
(put 'upcase-region 'disabled nil)
(put 'narrow-to-region 'disabled nil)
\end{verbatim}

\subsection{yes-or-no}
\label{sec:orgc4e9103}
Change all prompts to y or n:
\begin{verbatim}
(fset 'yes-or-no-p 'y-or-n-p)
\end{verbatim}
\subsection{shell env}
\label{sec:org08e86a9}
\begin{verbatim}
(use-package
  exec-path-from-shell
  :ensure
  :config
  (exec-path-from-shell-initialize))
\end{verbatim}
\subsection{always show line number?}
\label{sec:org8408fc3}
\begin{verbatim}
;;(global-display-line-numbers-mode)
\end{verbatim}
\section{Managing buffers}
\label{sec:orgd265ce4}
\subsection{winner-mode (layout managing)}
\label{sec:org85023ba}
Enables \texttt{winner-mode}. Navigate buffer-window configs with \texttt{C-c left}
and \texttt{C-c right}.

\begin{verbatim}
(winner-mode 1)
\end{verbatim}

\subsection{eyebrowse (layout managing)}
\label{sec:org8bb872b}
Awesome window manager. It's like using i3m but inside emacs. Use the
\texttt{C-c C-w <0..9>} key to switch to so called desktop. On each desktop,
you can have different buffers open and so on, so I don't have to
close buffers, or \texttt{C-x b} a lot anymore.

\begin{verbatim}
(use-package eyebrowse
  :ensure t)
\end{verbatim}

\subsection{toggle-maximize-buffer}
\label{sec:org16e5ee1}
Temporarily maximize a buffer.
\href{https://gist.github.com/mads379/3402786}{found here}

\begin{verbatim}
(defun toggle-maximize-buffer () "Maximize buffer"
  (interactive)
  (if (= 1 (length (window-list)))
      (jump-to-register '_)
    (progn
      (window-configuration-to-register '_)
      (delete-other-windows))))
\end{verbatim}

Map it to a key.
\begin{verbatim}
(global-set-key [M-f8] 'toggle-maximize-buffer)
\end{verbatim}
\subsection{select buffer}
\label{sec:orgc2a0896}
\begin{verbatim}
(use-package ace-window
  :ensure
  :config
  (setq aw-ignore-current t)
  (setq aw-keys '(?a ?s ?d ?f ?g ?h ?j ?k ?l))
  (setq aw-background nil)
  (global-set-key (kbd "M-o") 'ace-window)
  (custom-set-faces
   '(aw-leading-char-face
     ((t (:inherit ace-jump-face-foreground
                   :background "white"
                   :box nil)))))
)
\end{verbatim}
\section{Colors and Look}
\label{sec:org9d666d7}
\subsection{Theme}
\label{sec:org2786988}
Loading a theme I like:

\begin{verbatim}
  (use-package sublime-themes
    :ensure
    :config)
(add-hook 'after-init-hook (lambda () (load-theme 'spolsky t)))
\end{verbatim}
\subsection{Hide menu bar \& toolbar}
\label{sec:orgb3bf410}
Using i3 is forcing me to use keyboard.

\begin{verbatim}
(menu-bar-mode -1)
(toggle-scroll-bar -1)
(tool-bar-mode -1)
(blink-cursor-mode -1)
\end{verbatim}
\subsection{Be quiet}
\label{sec:org7cf1787}
Remove bell and dings.

\begin{verbatim}
(setq ring-bell-function
      '(lambda ()
         (message "The answer is 42...")))
(setq echo-keystrokes 0.1 use-dialog-box nil visible-bell t)
\end{verbatim}
\subsection{Fontlock}
\label{sec:org974af92}
This gets the font coloring switched on for all buffers.
\subsubsection{{\bfseries\sffamily TODO} Note: this should be the default, maybe this can go ?}
\label{sec:orgd00a88d}
\begin{verbatim}
(global-font-lock-mode t)
\end{verbatim}
\subsection{In terminal mode}
\label{sec:orge4083d5}
\begin{verbatim}
(when (display-graphic-p)
  (set-background-color "#ffffff")
  (set-foreground-color "#141312"))
\end{verbatim}
\subsection{In X11 mode: mouse and window title}
\label{sec:org07126f6}
\begin{verbatim}
(setq frame-title-format "emacs @ %b - %f")
(when window-system
  (mouse-wheel-mode)  ;; enable wheelmouse support by default
  (set-selection-coding-system 'compound-text-with-extensions))
\end{verbatim}
\subsection{Look: buffer naming}
\label{sec:org969cc56}
\begin{verbatim}
(use-package uniquify
  :init
  (setq uniquify-buffer-name-style 'post-forward-angle-brackets))
\end{verbatim}
\subsection{Buffer Decorations}
\label{sec:org5a3abbd}
Setup the visual cues about the current editing buffer
\begin{verbatim}
(column-number-mode t)
(setq visible-bell t)
(setq scroll-step 1)
(setq-default transient-mark-mode t)  ;; highlight selection
\end{verbatim}
\subsection{nyan-mode}
\label{sec:org95fe87b}
\begin{verbatim}
(use-package nyan-mode
  :ensure t
  :bind ("C-c f n" . 'nyan-mode))
\end{verbatim}
\subsection{dynamic cursor colors}
\label{sec:org61af5ce}
The cursor is displayed in different colors, depending on overwrite or insert mode.
\begin{verbatim}
(setq hcz-set-cursor-color-color "")
(setq hcz-set-cursor-color-buffer "")

(defun hcz-set-cursor-color-according-to-mode ()
  "change cursor color according to some minor modes."
  ;; set-cursor-color is somewhat costly, so we only call it when needed:
  (let ((color
         (if buffer-read-only "orange"
           (if overwrite-mode "red"
             "green"))))
    (unless (and
             (string= color hcz-set-cursor-color-color)
             (string= (buffer-name) hcz-set-cursor-color-buffer))
      (set-cursor-color (setq hcz-set-cursor-color-color color))
      (setq hcz-set-cursor-color-buffer (buffer-name)))))

(add-hook 'post-command-hook 'hcz-set-cursor-color-according-to-mode)
\end{verbatim}
\subsection{faces}
\label{sec:orgac16c39}
This makes some of the faces a bit more contrasted.
\begin{verbatim}
;; faces for general region highlighting zenburn is too low-key.
(custom-set-faces
 '(highlight ((t (:background "forest green"))))
 '(region ((t (:background "forest green")))))
\end{verbatim}
\textbf{*}
\subsection{delight}
\label{sec:org89e8bbb}
Package to remove some info from the mode-line for minor-modes.
\begin{verbatim}
(use-package delight
  :ensure t)
\end{verbatim}
\subsection{remove some modelines}
\label{sec:org341e96a}
\begin{verbatim}
(use-package eldoc
  :delight)
\end{verbatim}

\subsection{alternate key mappings}
\label{sec:org83c82f7}
Letting one enter chars that are otherwise difficult in e.g. the minibuffer.
\begin{verbatim}
(global-set-key (kbd "C-m") 'newline-and-indent)
(global-set-key (kbd "C-j") 'newline)
(global-set-key [delete] 'delete-char)
(global-set-key [kp-delete] 'delete-char)
\end{verbatim}
\subsection{Macros}
\label{sec:orgc4862c3}
\begin{verbatim}
(global-set-key [f3] 'start-kbd-macro)
(global-set-key [f4] 'end-kbd-macro)
(global-set-key [f5] 'call-last-kbd-macro)
\end{verbatim}
\subsection{Text size}
\label{sec:org80a7a91}
Increase/decrease text size
\begin{verbatim}
(define-key global-map (kbd "C-+") 'text-scale-increase)
(define-key global-map (kbd "C--") 'text-scale-decrease)
\end{verbatim}
\subsection{multiple regions}
\label{sec:orgd08da0d}
\begin{verbatim}
(global-set-key (kbd "C-M-i") 'iedit-mode)
\end{verbatim}
\subsection{Moving around buffers}
\label{sec:orgf32a8a4}
\begin{verbatim}
(global-set-key (kbd "C-c <C-left>")  'windmove-left)
(global-set-key (kbd "C-c <C-right>") 'windmove-right)
(global-set-key (kbd "C-c <C-up>")    'windmove-up)
(global-set-key (kbd "C-c <C-down>")  'windmove-down)
(global-set-key (kbd "C-c C-g") 'goto-line)
\end{verbatim}
\subsection{multiple-cursors}
\label{sec:org024f2db}
Configure the shortcuts for multiple cursors
\begin{verbatim}
(use-package multiple-cursors
  :ensure t
  :bind (("C-S-c C-S-c" . 'mc/edit-lines)
         ("C->" . 'mc/mark-next-like-this)
         ("C-<" . 'mc/mark-previous-like-this)
         ("C-c C->" . 'mc/mark-all-like-this)))
\end{verbatim}
\subsection{ace-jump-mode}
\label{sec:org94ba6a6}
Let's one jump around text
\begin{verbatim}
(use-package ace-jump-mode
  :ensure t
  :bind (("C-c C-SPC" . 'ace-jump-mode)
         ("C-c C-DEL" . 'ace-jump-mode-pop-mark)))
\end{verbatim}
\subsection{Hydra}
\label{sec:org73d59b8}
\begin{verbatim}
(use-package hydra
  :ensure t)
\end{verbatim}
\subsection{dimmer}
\label{sec:org82b1c2e}
\begin{verbatim}
(use-package dimmer
:ensure
:config
 (dimmer-configure-which-key)
 (dimmer-configure-helm))
 (dimmer-mode t)
\end{verbatim}
\subsection{Highlight indent}
\label{sec:org926d215}
\begin{verbatim}
(use-package highlight-indent-guides
:ensure
:config
(setq highlight-indent-guides-method 'character))
(add-hook 'prog-mode-hook 'highlight-indent-guides-mode)
\end{verbatim}
\section{Editing style}
\label{sec:org02e2826}
\subsection{No tabs, ever. No trailing spaces either.}
\label{sec:orgae78ce2}
\begin{verbatim}
(setq-default indent-tabs-mode nil)
(setq require-final-newline t)
(setq next-line-add-newlines nil)
(add-hook 'before-save-hook 'delete-trailing-whitespace)
\end{verbatim}
\subsection{Mark the 80 cols boundary}
\label{sec:org5332146}
\begin{verbatim}
(use-package column-enforce-mode
  :ensure t
  :config
  (setq column-enforce-column 80)
  :bind ("C-c m" . 'column-enforce-mode))
;; column-enforce-face
\end{verbatim}
\subsection{Better kill ring}
\label{sec:orgc1dbb31}
Don't use \texttt{popup-kill-ring} as it's dead. Use the \href{https://github.com/browse-kill-ring/browse-kill-ring}{browse-kill-ring}:

\begin{verbatim}
(use-package browse-kill-ring
:ensure
:config
(setq browse-kill-ring-highlight-current-entry t)
(setq browse-kill-ring-highlight-inserted-item t))

(browse-kill-ring-default-keybindings)
\end{verbatim}
\subsection{Parenthesis mode}
\label{sec:org99aaf71}
Borrowing from old init.el:
\begin{verbatim}
(show-paren-mode t)
(set-face-attribute 'region nil
                    :background "#666"
                    :foreground "#d52349")
(set-face-background 'show-paren-match (face-background 'default))
(set-face-foreground 'show-paren-match "#d52349")
(set-face-attribute 'show-paren-match nil
                    :weight 'extra-bold)
\end{verbatim}

\section{Cool packages}
\label{sec:orgb21806a}
\subsection{annotate-mode}
\label{sec:org2d43ee0}
The file-annotations are store externally. Seems to fail with \texttt{args-out-of-range} and then Emacs is confused. (filed issue for this)

Also, it seems to interfere with colorful modes like \texttt{magit} or \texttt{org-agenda-mode} so that I went with a whitelist instead of the wish of a blacklist of modes.

\begin{verbatim}
(use-package annotate
  :ensure t
  :bind ("C-c C-A" . 'annotate-annotate)  ;; for ledger-mode, as 'C-c C-a' is taken there.
  :config
  (add-hook 'org-mode 'annotate-mode)
  (add-hook 'csv-mode 'annotate-mode)
  (add-hook 'c-mode 'annotate-mode)
  (add-hook 'c++-mode 'annotate-mode)
  (add-hook 'sh-mode 'annotate-mode)
  (add-hook 'ledger-mode 'annotate-mode)
;;;  (define-globalized-minor-mode global-annotate-mode annotate-mode
;;;    (lambda () (annotate-mode 1)))
;;;  (global-annotate-mode 1)
  )
\end{verbatim}

\subsection{protobuf-mode}
\label{sec:orga1b4be7}
Mode for Google protocol buffer mode
\begin{verbatim}
(use-package protobuf-mode
  :ensure t
  :mode "\\.proto\\'")
\end{verbatim}
\subsection{Helm (list completion)}
\label{sec:org6fe00d9}
Trying out Helm instead of icicles, as it is available on ELPA.

I just took over the config described in this \href{https://tuhdo.github.io/helm-intro.html}{helm intro}.

\begin{verbatim}
(use-package helm
 :ensure t
 :delight helm-mode
 :config
  (require 'helm-config)
  ;; The default "C-x c" is quite close to "C-x C-c", which quits Emacs.
  ;; Changed to "C-c h". Note: We must set "C-c h" globally, because we
  ;; cannot change `helm-command-prefix-key' once `helm-config' is loaded.
  (global-set-key (kbd "C-c h") 'helm-command-prefix)
  (global-unset-key (kbd "C-x c"))

  (define-key helm-map (kbd "<tab>") 'helm-execute-persistent-action) ; rebind tab to run persistent action
  (define-key helm-map (kbd "C-c f") 'helm-execute-persistent-action) ; make TAB work in terminal
  (define-key helm-map (kbd "C-z")  'helm-select-action) ; list actions using C-z

  (when (executable-find "curl")
    (setq helm-google-suggest-use-curl-p t))

  (setq helm-split-window-inside-p            t ; open helm buffer inside current window, not occupy whole other window
        helm-move-to-line-cycle-in-source     t ; move to end or beginning of source when reaching top or bottom of source.
        helm-ff-search-library-in-sexp        t ; search for library in `require' and `declare-function' sexp.
        helm-scroll-amount                    8 ; scroll 8 lines other window using M-<next>/M-<prior>
        helm-ff-file-name-history-use-recentf t
        helm-echo-input-in-header-line t)

  (setq helm-autoresize-max-height 50)
  (setq helm-autoresize-min-height 0)
  (helm-autoresize-mode 1)

  (setq helm-M-x-fuzzy-match t)
  (setq helm-buffers-fuzzy-matching t
        helm-recentf-fuzzy-match    t)
  (setq helm-semantic-fuzzy-match t
        helm-imenu-fuzzy-match    t)
(setq helm-locate-fuzzy-match t)
(setq helm-apropos-fuzzy-match t)
(setq helm-lisp-fuzzy-completion t)
(helm-mode 1)

  (global-set-key (kbd "M-x") 'helm-M-x))
\end{verbatim}

I don't know what this is about, copied from the intro:

\begin{verbatim}
(defun spacemacs//helm-hide-minibuffer-maybe ()
  "Hide minibuffer in Helm session if we use the header line as input field."
  (when (with-helm-buffer helm-echo-input-in-header-line)
    (let ((ov (make-overlay (point-min) (point-max) nil nil t)))
      (overlay-put ov 'window (selected-window))
      (overlay-put ov 'face
                   (let ((bg-color (face-background 'default nil)))
                     `(:background ,bg-color :foreground ,bg-color)))
      (setq-local cursor-type nil))))

(add-hook 'helm-minibuffer-set-up-hook
          'spacemacs//helm-hide-minibuffer-maybe)
\end{verbatim}

Found \href{https://www.reddit.com/r/emacs/comments/30yer0/helm\_and\_recentf\_tips/}{this reddit post} of using \texttt{helm-mini}:
\begin{verbatim}
(setq helm-mini-default-sources '(helm-source-buffers-list
                                  helm-source-recentf
                                  helm-source-bookmarks
                                  helm-source-buffer-not-found))
\end{verbatim}

\subsection{\href{https://github.com/smihica/emmet-mode}{emmet-mode}}
\label{sec:orga84ad0e}
Useful abbreviations when coding in HTML.
\begin{verbatim}
(use-package emmet-mode
:ensure t)
\end{verbatim}
\subsection{rainbow-mode}
\label{sec:orgd6de9f2}
Colorize color names and codes in the correct color.
\begin{verbatim}
(use-package rainbow-mode
  :ensure t)
\end{verbatim}
\subsection{writeroom-mode}
\label{sec:org7c2a711}
\begin{verbatim}
(use-package writeroom-mode
  :ensure t
  :init
  (global-set-key (kbd "C-c f w") 'writeroom-mode))
\end{verbatim}

\subsection{wgrep-mode}
\label{sec:org5845f25}
\begin{verbatim}
(use-package wgrep
  :ensure t)
\end{verbatim}

\subsection{\href{https://github.com/ledger/ledger-mode}{ledger-mode}}
\label{sec:org6173c19}
\subsubsection{Cleanup ledger file}
\label{sec:orgfa61e63}
\begin{verbatim}
(defun single-lines-only ()
  "replace multiple blank lines with a single one"
  (interactive)
  (goto-char (point-min))
  (while (re-search-forward "\\(^\\s-*$\\)\n" nil t)
    (replace-match "\n")
    (forward-char 1)))

(defun paf/cleanup-ledger-buffer ()
  "Cleanup the ledger file"
  (interactive)
  (delete-trailing-whitespace)
  (single-lines-only)
  (ledger-mode-clean-buffer)
  (ledger-sort-buffer))
\end{verbatim}
\subsubsection{Setup}
\label{sec:orgcd35993}
\begin{verbatim}
(use-package ledger-mode
  :ensure nil
  :pin manual
  :mode "\\.ledger\\'"
  :bind ("<f6>" . 'paf/cleanup-ledger-buffer)
  :config
  (setq ledger-reconcile-default-commodity "CHF"))
\end{verbatim}
\subsection{\href{http://www.gnu.org/software/hyperbole/}{hyperbole}}
\label{sec:org8b68213}
Let's try this too, even though I do not quite get the point of this
whole package yet.

\textbf{NOTE} assigns \texttt{hui-search-web} to \texttt{C-c C-/} to not clobber the later used \texttt{C-c /} from OrgMode (org-mode sparse trees). This works because hyperbole will first check if the function is already bound to some key before binding it to the coded default.
\begin{verbatim}
(use-package hyperbole
  :ensure t
  :config
  (bind-key "C-c C-/" 'hui-search-web)  ;; bind before calling require
  (require 'hyperbole))
\end{verbatim}
\subsection{\href{https://github.com/fourier/ztree\#ztree}{ztree}}
\label{sec:orga0f31e5}
A tree-view navigation of files, with diff tool for directories.
\begin{verbatim}
(use-package ztree
  :ensure t)
\end{verbatim}
\subsection{anzu}
\label{sec:org09540b2}
Show number of search matches.
\begin{verbatim}
(use-package
  anzu

  :ensure
  :config)
(global-anzu-mode +1)
\end{verbatim}
\subsection{whitespace}
\label{sec:orga125f16}
Highlight unnecessary chars and lines over 80.
\begin{verbatim}
(use-package
  whitespace

  :ensure
  :config (setq whitespace-style '(face empty tabs lines-tail trailing))
  :config (global-whitespace-mode t))
\end{verbatim}
\subsection{move-text}
\label{sec:org09124e8}
Looks convenient:
\begin{verbatim}
  (use-package move-text
    :ensure
    :config)
(move-text-default-bindings)
\end{verbatim}
\subsection{string inflection}
\label{sec:orgadb1cdf}
Change text style quickly:
\begin{verbatim}
(use-package string-inflection
  :ensure
  :config)
(add-hook 'elpy-mode-hook
            '(lambda ()
               (local-set-key (kbd "C-q C-u") 'string-inflection-python-style-cycle)))
(add-hook 'org-mode-hook
            '(lambda ()
               (local-set-key (kbd "C-q C-u") 'string-inflection-python-style-cycle)))
(add-hook 'mu4e-compose-mode-hook
            '(lambda ()
               (local-set-key (kbd "C-q C-u") 'string-inflection-python-style-cycle)))
\end{verbatim}
\subsection{guide-key}
\label{sec:orgb8b5986}
Show key bindings. The built-in one is \texttt{C-h b}.
\begin{verbatim}
(use-package guide-key
  :ensure
  :config
  (setq guide-key/guide-key-sequence t)
  (setq guide-key/highlight-command-regexp "rectangle"))
(guide-key-mode 1)
\end{verbatim}


\section{Coding}
\label{sec:orgea3ed7c}
\subsection{Version control}
\label{sec:orgaa90700}
\subsubsection{magit}
\label{sec:org69f0d1d}
Add the powerful Magit
\begin{verbatim}
(use-package magit
  :ensure t
  :defer
  :config
  (add-hook 'after-save-hook 'magit-after-save-refresh-status t)
  :bind ("C-x g" . 'magit-status))
(use-package magit-todos
  :ensure t
  :defer)
(use-package
  magit-gitflow

  :ensure
  :config (add-hook 'magit-mode-hook 'turn-on-magit-gitflow))
\end{verbatim}
\subsubsection{monky}
\label{sec:org9f1e752}
Add the Magit-copy for Mercurial 'monky'
\begin{verbatim}
(use-package monky
  :ensure t
  :defer
  :bind ("C-x m" . 'monky-status))
\end{verbatim}
\subsubsection{Global caller}
\label{sec:org8455c22}
Have a single binding to call the most appropriate tool given the repository.
\begin{verbatim}
(defun paf/vcs-status ()
     (interactive)
     (condition-case nil
         (magit-status-setup-buffer)
       (error (monky-status))))

(global-set-key (kbd "C-c f v") 'paf/vcs-status)
\end{verbatim}

\subsection{Projectile}
\label{sec:org15340ac}
Start using projectile. It has the documentation \href{https://docs.projectile.mx/en/latest/}{here}.
\begin{verbatim}
(use-package projectile
  :ensure t
  :config
  (define-key projectile-mode-map (kbd "C-c p") 'projectile-command-map)
  (setq projectile-completion-system 'helm)
  (projectile-mode +1))

(use-package helm-projectile
  :ensure t
  :after projectile
  :requires projectile
  :delight projectile-mode
  :config
  (helm-projectile-on))
\end{verbatim}

Also make sure we do have the faster \href{https://github.com/ggreer/the\_silver\_searcher\#the-silver-searcher}{silver searcher} version.  This
may need you to install the corresponding tool for this, with the
following snippet:
\begin{verbatim}
# helm-ag uses this for faster grepping
if [[ "$(uname)" == "Darwin" ]]; then
  install_pkg -x ag the_silver_searcher
else
  install_pkg -x ag silversearcher-ag
fi
\end{verbatim}

Search the entire project with \texttt{C-c p s s} for a regexp. This let's
you turn the matching results into an editable buffer using \texttt{C-c
C-e}. Other keys are listed \href{https://github.com/syohex/emacs-helm-ag\#keymap}{here}.

\begin{verbatim}
(use-package helm-ag
  :ensure t
  :config)
\end{verbatim}

I havae used it by a \texttt{M-?} binding. It's just old habit:
\begin{verbatim}
(global-set-key (kbd "M-?") 'helm-ag)
\end{verbatim}

\subsection{no indentation of namespaces in C++}
\label{sec:org95d7e4d}
Essentially, use the Google C++ style formatting.
\begin{verbatim}
(use-package google-c-style
  :ensure t
  :config
  (add-hook 'c-mode-common-hook 'google-set-c-style)
  (add-hook 'c-mode-common-hook 'google-make-newline-indent))

(use-package flymake-google-cpplint
  :ensure t)
\end{verbatim}
\subsection{search \& jump}
\label{sec:org25c8cd2}
\subsubsection{ag}
\label{sec:orgc166b4c}
Use the silversearcher.
\begin{verbatim}
(use-package ag
  :ensure
  :config
  (setq ag-highlight-search t)
  (setq ag-reuse-buffers t))
\end{verbatim}

Follow the \href{https://github.com/emacsorphanage/helm-ag}{helm-ag manual}, "Insert thing at point as default search
pattern, if this value is non nil":

\begin{verbatim}
(setq helm-ag-insert-at-point 'symbol)
(setq helm-ag-use-temp-buffer t)
\end{verbatim}
\subsubsection{dumb-jump}
\label{sec:org223f303}

First, let's make sure we have \texttt{xref} because we will hook into the
xref backend:
\begin{verbatim}
(use-package xref
  :ensure
  :config)
\end{verbatim}

Now install \texttt{dumb-jump}:
\begin{verbatim}
(use-package dumb-jump
  :ensure
  :config
  (setq dumb-jump-prefer-searcher 'ag))
\end{verbatim}

Last, hook to \texttt{xref} to use \texttt{M.} bind:
\begin{verbatim}
(add-hook 'xref-backend-functions #'dumb-jump-xref-activate)
\end{verbatim}
\subsubsection{ripgrep}
\label{sec:orgc93363d}
This enables searching recursively in projects.
\begin{verbatim}
# This can be used by helm-ag for faster grepping
install_pkg -x rg ripgrep
\end{verbatim}

\begin{verbatim}
(use-package ripgrep
  :ensure t)
(use-package projectile-ripgrep
  :ensure t
  :requires (ripgrep projectile))
\end{verbatim}

\subsection{commenting out}
\label{sec:orge08d88e}
Easy commenting out of lines.
\begin{verbatim}
(autoload 'comment-out-region "comment" nil t)
(global-set-key (kbd "C-c q") 'comment-out-region)
\end{verbatim}

\subsection{Deduplicate and sort}
\label{sec:orgf63ef60}
Help cleanup the includes and using lists.
\href{http://www.emacswiki.org/emacs/DuplicateLines}{found here}
\begin{verbatim}
(defun uniquify-region-lines (beg end)
  "Remove duplicate adjacent lines in region."
  (interactive "*r")
  (save-excursion
    (goto-char beg)
    (while (re-search-forward "^\\(.*\n\\)\\1+" end t)
      (replace-match "\\1"))))

(defun paf/sort-and-uniquify-region ()
  "Remove duplicates and sort lines in region."
  (interactive)
  (sort-lines nil (region-beginning) (region-end))
  (uniquify-region-lines (region-beginning) (region-end)))
\end{verbatim}

Simplify cleanup of \texttt{\#include} / \texttt{typedef} / \texttt{using} blocks.
\begin{verbatim}
(global-set-key (kbd "C-c f s") 'paf/sort-and-uniquify-region)
\end{verbatim}

\subsection{diffing}
\label{sec:orge8d1f8d}
\href{https://github.com/justbur/emacs-vdiff}{vdiff} let's one compare buffers or files.
\begin{verbatim}
(use-package vdiff
  :ensure t
  :config
  ; This binds commands under the prefix when vdiff is active.
  (define-key vdiff-mode-map (kbd "C-c") vdiff-mode-prefix-map))
\end{verbatim}

\subsection{yasnippet}
\label{sec:orgee60fa9}
Let's first see how far I get with file-based capture templates and yankpad.
\begin{verbatim}
  (use-package yasnippet
    :ensure t)
  (use-package auto-yasnippet
    :ensure t
:after yasnippet
    :config
    (bind-key "C-c f C-s c" 'aya-create)
    (bind-key "C-c f C-s e" 'aya-expand))
\end{verbatim}

\subsection{Selective display}
\label{sec:orga61c5d6}
Will fold all text indented more than the position of the cursor at
the time the keys are pressed.

\begin{verbatim}
(defun set-selective-display-dlw (&optional level)
  "Fold text indented more than the cursor.
   If level is set, set the indent level to level.
   0 displays the entire buffer."
  (interactive "P")
  (set-selective-display (or level (current-column))))

(global-set-key "\C-x$" 'set-selective-display-dlw)
\end{verbatim}
\subsection{Info in the gutter}
\label{sec:org361b679}
\subsubsection{git informations}
\label{sec:orged59afe}
\begin{verbatim}
(use-package git-gutter-fringe+
  :ensure t
  :defer
  :if window-system
  :bind ("C-c g" . 'git-gutter+-mode))
\end{verbatim}
\subsection{Speedup VCS}
\label{sec:orga358b69}
Regexp matching directory names that are not under VC's control. The
default regexp prevents fruitless and time-consuming attempts to
determine the VC status in directories in which filenames are
interpreted as hostnames.

\begin{verbatim}
(defvar locate-dominating-stop-dir-regexp
  "\\`\\(?:[\\/][\\/][^\\/]+\\|/\\(?:net\\|afs\\|\\.\\.\\.\\)/\\)\\'")
\end{verbatim}
\subsection{Dealing with numbers}
\label{sec:orgb1d1489}
Simple way to increase/decrease a number in code.
\begin{verbatim}
(use-package shift-number
  :ensure t
  :bind (("M-+" . shift-number-up)
         ("M-_" . shift-number-down)))
\end{verbatim}
\subsection{GDB with many windows}
\label{sec:orge23bb48}
\subsubsection{{\bfseries\sffamily TODO} Make it so that the source frame placement is forced only when using gdb.}
\label{sec:orgd254f54}

\begin{verbatim}
(setq gdb-many-windows t)
(setq gdb-use-separate-io-buffer t)

(defun easy-gdb-top-of-stack-and-restore-windows ()
  (interactive)
  (switch-to-buffer (gdb-stack-buffer-name))
  (goto-char (point-min))
  (gdb-select-frame)
  (gdb-restore-windows)
  (other-window 2))

(global-set-key (kbd "C-x C-a C-t") 'easy-gdb-top-of-stack-and-restore-windows)
\end{verbatim}

This should display the source code always in the same window when debugging.
Found on \href{https://stackoverflow.com/questions/39762833/emacsgdb-customization-how-to-display-source-buffer-in-one-window}{Stack Overflow}.
\begin{verbatim}
; This unfortunately also messes up the regular frame navigation of source code.
;(add-to-list 'display-buffer-alist
;             (cons 'cdb-source-code-buffer-p
;                   (cons 'display-source-code-buffer nil)))

(defun cdb-source-code-buffer-p (bufName action)
  "Return whether BUFNAME is a source code buffer."
  (let ((buf (get-buffer bufName)))
    (and buf
         (with-current-buffer buf
           (derived-mode-p buf 'c++-mode 'c-mode 'csharp-mode 'nxml-mode)))))

(defun display-source-code-buffer (sourceBuf alist)
  "Find a window with source code and set sourceBuf inside it."
  (let* ((curbuf (current-buffer))
         (wincurbuf (get-buffer-window curbuf))
         (win (if (and wincurbuf
                       (derived-mode-p sourceBuf 'c++-mode 'c-mode 'nxml-mode)
                       (derived-mode-p (current-buffer) 'c++-mode 'c-mode 'nxml-mode))
                  wincurbuf
                (get-window-with-predicate
                 (lambda (window)
                   (let ((bufName (buffer-name (window-buffer window))))
                     (or (cdb-source-code-buffer-p bufName nil)
                         (assoc bufName display-buffer-alist)
                         ))))))) ;; derived-mode-p doesn't work inside this, don't know why...
    (set-window-buffer win sourceBuf)
    win))
\end{verbatim}

Here is my cheatsheet for the keyboard commands:

All prefixed with \texttt{C-x C-a}

\begin{center}
\begin{tabular}{llc}
\hline
Domain & Command & C-<key>\\
\hline
Breakpoint & set & b\\
 & temporary & t\\
 & delete & d\\
\hline
Execute & Next & n\\
 & Step Into & s\\
 & Return / Finish & f\\
 & Continue (run) & r\\
\hline
Stack & Up & <\\
 & Down & >\\
\hline
Execute & Until current line & u\\
(rarer) & Single instruction & i\\
 & Jump to current line & j\\
\hline
\end{tabular}
\end{center}

\subsection{smartparens}
\label{sec:org345e6f2}
Well, who wants to type parenthesis:
\begin{verbatim}
(use-package smartparens
  :ensure
  :config
  (require 'smartparens-config))
(add-hook 'prog-mode-hook #'smartparens-mode)
\end{verbatim}
\subsection{specific languages}
\label{sec:orgf435760}
These are minor modes to handle programming language specifics which
are often termed as development rules agreed by the team.
\subsubsection{c}
\label{sec:org699daa1}
\begin{enumerate}
\item header/implementation toggle
\label{sec:orgafc8fc3}

Switch from header to implementation file quickly.
\begin{verbatim}
(add-hook 'c-mode-common-hook
          (lambda ()
            (local-set-key  (kbd "C-c o") 'ff-find-other-file)))
\end{verbatim}
\end{enumerate}
\subsubsection{web-mode}
\label{sec:orgd193111}
web-mode with config for Polymer editing
\begin{verbatim}
(use-package web-mode
  :ensure t
  :mode "\\.html\\'"
  :config
  (setq web-mode-markup-indent-offset 2)
  (setq web-mode-css-indent-offset 2)
  (setq web-mode-code-indent-offset 2)
  (setq web-mode-enable-current-element-highlight t)
  (setq web-mode-enable-current-column-highlight t)
  (setq web-mode-enable-css-colorization t)
  (setq web-mode-comment-style 2))
\end{verbatim}
\subsubsection{csv-mode}
\label{sec:org2d6235a}
mode to edit CSV files.
\begin{verbatim}
(use-package csv-mode
  :ensure t
  :mode "\\.csv\\'")
\end{verbatim}
\subsubsection{taskjuggler-mode (tj3-mode)}
\label{sec:org801e7d4}
\begin{verbatim}
(use-package tj3-mode
  :ensure t
  :after org-plus-contrib
  :config
  (require 'ox-taskjuggler)
  (custom-set-variables
   '(org-taskjuggler-process-command "/usr/local/bin/tj3 --silent --no-color --output-dir %o %f")
   '(org-taskjuggler-project-tag "PRJ")))
\end{verbatim}
\subsubsection{markdown-mode}
\label{sec:org188de1d}
Enough to handle my Markdown needs.

\begin{verbatim}
(use-package writegood-mode
  :ensure
  :config)

(use-package
  markdown-mode

  :ensure
  :config
  (add-to-list 'auto-mode-alist '("\\.md\\'" . markdown-mode))
  (add-hook 'markdown-mode-hook
            (lambda ()
              (visual-line-mode t)
              (writegood-mode t)
              (flyspell-mode t))))
\end{verbatim}
\subsubsection{python}
\label{sec:org8e9ecf7}

Setup an IDE:
\begin{verbatim}
(use-package elpy
  :ensure
  :init
  :init
  (elpy-enable))
\end{verbatim}

ELPY has its own indentation mode, which is overriding the one I use globally, so disable this one:
\begin{verbatim}
(add-hook 'elpy-mode-hook (lambda () (highlight-indentation-mode -1)))
\end{verbatim}

\begin{enumerate}
\item py-autopep8
\label{sec:org22160f0}
Add hook to reformat python code based on pep8 spec. You need to
install \texttt{pip install autopep8} offline.

\begin{verbatim}
(use-package
  py-autopep8

  :ensure t
  :config
  (add-hook 'python-mode-hook 'py-autopep8-enable-on-save)
  (setq py-autopep8-options '("--max-line-length=80")))
\end{verbatim}
\item py-isort
\label{sec:org8907059}
Sort python import. Need to install \texttt{pip install isort} offline.

\begin{verbatim}
(use-package
  py-isort

  :ensure
  :config
  (add-hook 'before-save-hook 'py-isort-before-save)
  (setq py-isort-options '("-sl")))
\end{verbatim}
\item smartparens
\label{sec:org7655575}
\begin{verbatim}
(add-hook 'python-mode-hook #'smartparens-mode)
\end{verbatim}
\end{enumerate}
\subsubsection{javascript family: .js .ts .jsx}
\label{sec:org22e2d18}
There are a couple packages for .js files.

\begin{enumerate}
\item js2-mode
\label{sec:orgd4743b6}
First, use \texttt{js2-mode} to handle \texttt{.js} and \texttt{.jsx} files.

\begin{verbatim}
(use-package js2-mode
  :ensure
  :config
  (add-to-list 'auto-mode-alist '("\\.js[x]\\'" . js2-mode))
  (add-to-list 'auto-mode-alist '("\\.ts\\'" . js2-mode)))
(add-hook 'js-mode-hook #'smartparens-mode)
\end{verbatim}

\item prettier-js
\label{sec:orgc68b506}
Link js2-mode to prettier to beautify my code.
\begin{verbatim}
(use-package prettier-js
  :ensure
  :config
  (setq prettier-js-args '(
                           "--trailing-comma" "es5"
                           "--single-quote" "false"
                           "--print-width" "80"
                           "--tab-width" "2"
                           )));
(add-hook 'js2-mode-hook 'prettier-js-mode)
\end{verbatim}
\item js-doc
\label{sec:orge2df120}
Nothing is complete without a doc solution.

\begin{verbatim}
(use-package js-doc
  :ensure
  :config
  (setq js-doc-mail-address "fxia1@lenovo.com")
  (setq js-doc-author (format "Feng Xia <%s>" js-doc-mail-address))
  (setq js-doc-url "http://www.lenovo.com")
  (setq js-doc-license "Lenovo License")
)
(add-hook 'js2-mode-hook
          #'(lambda ()
              (define-key js2-mode-map "\C-ci" 'js-doc-insert-function-doc)
              (define-key js2-mode-map "@" 'js-doc-insert-tag)))

\end{verbatim}
\end{enumerate}
\section{Org}
\label{sec:org7b51ad7}
Load all my org stuff, but first org-mode itself.
\subsection{Init}
\label{sec:org596c0d2}
If variable \texttt{org-directory} is not set yet, map it to my home's
files. You may set this in the \texttt{\textasciitilde{}/.emacs} to another value,
e.g. \texttt{(setq org-directory "/ssh:fleury@machine.site.com:OrgFiles")}

\subsubsection{{\bfseries\sffamily NEXT} This does not seem to work, check out doc about \href{https://stackoverflow.com/questions/3806423/how-can-i-get-a-variables-initial-value-in-elisp}{defcustom}}
\label{sec:orgacdd28d}
\begin{verbatim}
(use-package org
  :ensure nil
  :delight org-mode
  :config
  (setq org-startup-indented t)
  (if (not (boundp 'org-directory))
      (setq org-directory "~/org"))
  (add-hook 'org-mode-hook #'(lambda ()
                               (visual-line-mode)
                               (org-indent-mode))))
\end{verbatim}
\subsection{Access my remote Org files directory}
\label{sec:org34ced97}
Let's bind this to a key, so I can open remote dirs. I suually put this in my \texttt{.emacs} as it is host- and user-specific.
\begin{verbatim}
(defun paf/open-remote-org-dir ()
  (interactive)
  (dired "/ssh:remote.host.com:org"))

(global-set-key (kbd "C-c f r o") 'paf/open-remote-org-dir)
\end{verbatim}
\subsection{Helper Functions / Tools found on the web / worg}
\label{sec:orgd4a35f1}
\subsubsection{Open remote org dir}
\label{sec:org7d25c24}
In your \texttt{.emacs} just add this to configure the location:

\begin{verbatim}
(setq remote-org-directory "/ssh:fleury@my.hostname.com:OrgFiles")
\end{verbatim}

Then you can use the keyboard shortcut to open that dir.

\begin{verbatim}
(defcustom remote-org-directory "~/OrgFiles"
  "Location of remove OrgFile directory, should you have one."
  :type 'string
  :group 'paf)
(defun paf/open-remote-org-directory ()
  (interactive)
  (find-file remote-org-directory))

(global-set-key (kbd "C-c f r o") 'paf/open-remote-org-directory)
\end{verbatim}

\subsubsection{org-protocol}
\label{sec:orgf4035e6}
Let other tools use emacs client to interact
\begin{verbatim}
(require 'org-protocol)
\end{verbatim}
\subsubsection{Org-relative file function}
\label{sec:orgac3a293}
\begin{verbatim}
(defun org-relative-file (filename)
  "Compute an expanded absolute file path for org files"
  (expand-file-name filename org-directory))
\end{verbatim}
\subsubsection{Adjust tags on the right}
\label{sec:orgb665931}
Dynamically adjust tag position
\href{https://orgmode.org/worg/org-hacks.html\#org0560357}{source on worg}

\begin{verbatim}
(defun ba/org-adjust-tags-column-reset-tags ()
  "In org-mode buffers it will reset tag position according to
`org-tags-column'."
  (when (and
         (not (string= (buffer-name) "*Remember*"))
         (eql major-mode 'org-mode))
    (let ((b-m-p (buffer-modified-p)))
      (condition-case nil
          (save-excursion
            (goto-char (point-min))
            (command-execute 'outline-next-visible-heading)
            ;; disable (message) that org-set-tags generates
            (cl-letf (((symbol-function 'message) #'format))
              (org-set-tags 1 t))
            (set-buffer-modified-p b-m-p))
        (error nil)))))

(defun ba/org-adjust-tags-column-now ()
  "Right-adjust `org-tags-column' value, then reset tag position."
  (set (make-local-variable 'org-tags-column)
       (- (- (window-width) (length org-ellipsis))))
  (ba/org-adjust-tags-column-reset-tags))

(defun ba/org-adjust-tags-column-maybe ()
  "If `ba/org-adjust-tags-column' is set to non-nil, adjust tags."
  (when ba/org-adjust-tags-column
    (ba/org-adjust-tags-column-now)))

(defun ba/org-adjust-tags-column-before-save ()
  "Tags need to be left-adjusted when saving."
  (when ba/org-adjust-tags-column
     (setq org-tags-column 1)
     (ba/org-adjust-tags-column-reset-tags)))

(defun ba/org-adjust-tags-column-after-save ()
  "Revert left-adjusted tag position done by before-save hook."
  (ba/org-adjust-tags-column-maybe)
  (set-buffer-modified-p nil))

;; between invoking org-refile and displaying the prompt (which
;; triggers window-configuration-change-hook) tags might adjust,
;; which invalidates the org-refile cache
(defadvice org-refile (around org-refile-disable-adjust-tags)
  "Disable dynamically adjusting tags"
  (let ((ba/org-adjust-tags-column nil))
    ad-do-it))
(ad-activate 'org-refile)

;; Now set it up
(setq ba/org-adjust-tags-column t)
;; automatically align tags on right-hand side
;; TODO(fleury): Does not seem to work as of 2017/12/18
;; Seems to work again 2018/11/01
(add-hook 'window-configuration-change-hook
          'ba/org-adjust-tags-column-maybe)
(add-hook 'before-save-hook 'ba/org-adjust-tags-column-before-save)
(add-hook 'after-save-hook 'ba/org-adjust-tags-column-after-save)
(add-hook 'org-agenda-mode-hook (lambda ()
                                  (setq org-agenda-tags-column (- (window-width)))))
\end{verbatim}

\begin{enumerate}
\item {\bfseries\sffamily TODO} Update \texttt{org-set-tags-to}
\label{sec:orge014cd6}
\href{https://orgmode.org/worg/doc.html\#org-set-tags-to}{\texttt{org-set-tags-to}} is gone, and \texttt{org-set-tags} with > 1 args is not working.
Not sure what to replace it with though\ldots{}
\end{enumerate}

\subsubsection{Preserve structure in archives}
\label{sec:orgbb3bb96}
Make sure archiving preserves the same tree structure, including when archiving subtrees.
\href{https://orgmode.org/worg/org-hacks.html\#org4265b4c}{source on worg}

\begin{verbatim}
(defun my-org-inherited-no-file-tags ()
  (let ((tags (org-entry-get nil "ALLTAGS" 'selective))
        (ltags (org-entry-get nil "TAGS")))
    (mapc (lambda (tag)
            (setq tags
                  (replace-regexp-in-string (concat tag ":") "" tags)))
          (append org-file-tags (when ltags (split-string ltags ":" t))))
    (if (string= ":" tags) nil tags)))
\end{verbatim}

This used to work, but \texttt{org-extract-archive-file} is no longer defined.
\begin{verbatim}
(defadvice org-archive-subtree
    (around my-org-archive-subtree-low-level activate)
  (let ((tags (my-org-inherited-no-file-tags))
        (org-archive-location
         (if (save-excursion (org-back-to-heading)
                             (> (org-outline-level) 1))
             (concat (car (split-string org-archive-location "::"))
                     "::* "
                     (car (org-get-outline-path)))
           org-archive-location)))
    ad-do-it
    (with-current-buffer (find-file-noselect (org-extract-archive-file))
      (save-excursion
        (while (org-up-heading-safe))
        (org-set-tags tags)))))
\end{verbatim}
\subsubsection{Auto-Refresh Agenda}
\label{sec:org63e5f41}
Refresh org-mode agenda regularly.
\href{https://orgmode.org/worg/org-hacks.html\#orgab827a7}{source on worg}
There are two functions that supposedly do the same.
\begin{verbatim}
(defun kiwon/org-agenda-redo-in-other-window ()
  "Call org-agenda-redo function even in the non-agenda buffer."
  (interactive)
  (let ((agenda-window (get-buffer-window org-agenda-buffer-name t)))
    (when agenda-window
      (with-selected-window agenda-window (org-agenda-redo)))))

(defun update-agenda-if-visible ()
  (interactive)
  (let ((buf (get-buffer "*Org Agenda*"))
        wind)
    (if buf
        (org-agenda-redo))))
\end{verbatim}
\subsubsection{Display Agenda when idle}
\label{sec:org5dd2e6d}
Show the agenda when emacs left idle.
\href{https://orgmode.org/worg/org-hacks.html\#orgaea636d}{source on worg}
\begin{verbatim}
(defun jump-to-org-agenda ()
  (interactive)
  (let ((buf (get-buffer "*Org Agenda*"))
        wind)
    (if buf
        (if (setq wind (get-buffer-window buf))
            (select-window wind)
          (if (called-interactively-p 'any)
              (progn
                (select-window (display-buffer buf t t))
                (org-fit-window-to-buffer)
                (org-agenda-redo)
                )
            (with-selected-window (display-buffer buf)
              (org-fit-window-to-buffer)
              ;;(org-agenda-redo)
              )))
      (call-interactively 'org-agenda-list)))
  ;;(let ((buf (get-buffer "*Calendar*")))
  ;;  (unless (get-buffer-window buf)
  ;;    (org-agenda-goto-calendar)))
  )
\end{verbatim}
\subsubsection{Display location in agenda}
\label{sec:org9f75eed}
From some help on \href{https://emacs.stackexchange.com/questions/26249/customize-text-after-task-in-custom-org-agenda-view}{this page} I think this could work:
\begin{verbatim}
(defun paf/org-agenda-get-location()
  "Gets the value of the LOCATION property"
  (let ((loc (org-entry-get (point) "LOCATION")))
    (if (> (length loc) 0)
        loc
      "")))
\end{verbatim}

Also, to set this after org-mode has loaded (\href{https://emacs.stackexchange.com/questions/19091/how-to-set-org-agenda-prefix-format-before-org-agenda-starts}{see here}):
\begin{verbatim}
(with-eval-after-load 'org-agenda
  (add-to-list 'org-agenda-prefix-format
               '(agenda . "  %-12:c%?-12t %(paf/org-agenda-get-location)% s"))
\end{verbatim}
\subsubsection{org-gtasks}
\label{sec:org36b9a30}
Should follow this git repo: \href{https://github.com/JulienMasson/org-gtasks}{org-gtasks}
I have copied a version of the file here, it's not yet available on MELPA.

To help debug, use this before running things:
(setq request-log-level 'debug)

\begin{verbatim}
(use-package request
  :ensure t)
(load-file "~/Emacs/org-gtasks.el")
\end{verbatim}

I have this currently in my `\textasciitilde{}/.emacs`:
\begin{verbatim}
(use-package org-gtasks
  :init
  (org-gtasks-register-account
     :name "pascal"
     :directory "~/OrgFiles/GTasks/"
     :client-id "XXX"
     :client-secret "XXX"))
\end{verbatim}
\subsubsection{org-super-agenda}
\label{sec:org71187a3}
This enables a more fine-grained filtering of the agenda items.
\begin{verbatim}
(use-package org-super-agenda
  :ensure t
  :config
  (org-super-agenda-mode t))
\end{verbatim}
\subsubsection{org-roam}
\label{sec:orgc5233bc}
My cheat sheet for \texttt{org-roam}

All keys prefixed with \texttt{C-c n}

\begin{center}
\begin{tabular}{lc}
\hline
Function & \texttt{C-c n <key>}\\
\hline
Toggle side panel & l\\
\hline
Find/create & f\\
Insert link & i\\
Capture & c\\
\hline
Graph & g\\
Switch to buffer & b\\
\hline
\end{tabular}
\end{center}


\begin{verbatim}
(use-package org-roam
  :ensure t
  :hook (after-init . org-roam-mode)
  :init (setq org-roam-directory
              (org-relative-file "OrgRoam"))
  :bind (:map org-roam-mode-map
              (("C-c n l" . org-roam)
               ("C-c n b" . org-roam-switch-to-buffer)
               ("C-c n f" . org-roam-find-file)
               ("C-c n c" . org-roam-capture)
               ("C-c n g" . org-roam-graph))
         :map org-mode-map
              (("C-c n i" . org-roam-insert))))

(use-package company-org-roam
  :ensure t
  :after org-roam)
\end{verbatim}

EmacSQL will need to get its C-binary compiled, and needs supporting tools. Note that 'tcc' for Termux seems not complete enough for the job.
\begin{verbatim}
# org-roam needs this binary
install_pkg -x sqlite3 sqlite3
# Make sure there is a C compiler for emacsql-sqlite
[[ -n "$(which cc)" ]] || install_pkg -x cc clang
\end{verbatim}

\subsubsection{org-clock-convenience}
\label{sec:orga9fd33e}
\begin{verbatim}
(use-package org-clock-convenience
  :ensure t
  :bind (:map org-agenda-mode-map
           ("<S-right>" . org-clock-convenience-timestamp-up)
           ("<S-left>" . org-clock-convenience-timestamp-down)
           ("[" . org-clock-convenience-fill-gap)
           ("]" . org-clock-convenience-fill-gap-both)))
\end{verbatim}
\subsubsection{org-kanban}
\label{sec:org06bed54}
\begin{verbatim}
;;(use-package org-kanban
;;  :ensure t)
\end{verbatim}
\subsubsection{org-board}
\label{sec:org4d0a7b7}
Archive entire sites locally with `wget`.
\begin{verbatim}
(use-package org-board
  :ensure t
  :config
  (global-set-key (kbd "C-c o") org-board-keymap))
\end{verbatim}

This is the needed tool used to fetch a URL's content.
\begin{verbatim}
# wget used for org-board archiving.
install_pkg -x wget wget
\end{verbatim}

\subsubsection{org-reveal}
\label{sec:org1ead49f}
This presentation generator is still under review (by me).

\begin{verbatim}
# Install reveal.js
if [[ -d "${HOME}/workspace/me/moment/reveal.js" ]]; then
  echo "Reveal already installed"
else
  (cd ~/ && git clone https://github.com/hakimel/reveal.js.git)
fi
\end{verbatim}

\begin{verbatim}
(use-package ox-reveal
  :ensure t
  :after (htmlize)
  :config
  (setq org-reveal-root (expand-file-name "~/workspace/me/moment/reveal.js")))

(use-package htmlize
  :ensure t)
\end{verbatim}

\subsubsection{iimage (M-I)}
\label{sec:orgdc30902}
Make the display of images a simple key-stroke away.
\begin{verbatim}
(defun paf/org-toggle-iimage-in-org ()
  "display images in your org file"
  (interactive)
  (if (face-underline-p 'org-link)
      (set-face-underline 'org-link nil)
    (set-face-underline 'org-link t))
  (iimage-mode 'toggle))

(use-package iimage
  :config
  (add-to-list 'iimage-mode-image-regex-alist
               (cons (concat "\\[\\[file:\\(~?" iimage-mode-image-filename-regex
                             "\\)\\]")  1))
  (add-hook 'org-mode-hook (lambda ()
                             ;; display images
                             (local-set-key "\M-I" 'paf/org-toggle-iimage-in-org)
                            )))
\end{verbatim}
\subsubsection{Properties collector}
\label{sec:org79c615f}
Collect properties into tables. See documentation in the file.
\begin{verbatim}
(load-file "~/Emacs/org-collector.el")
\end{verbatim}

\subsection{My Setup}
\label{sec:org20c29c0}
These are mostly org-config specific to me, myself and I.
\subsubsection{key mappings}
\label{sec:org14cf40d}
\begin{verbatim}
(global-set-key (kbd "C-c l") 'org-store-link)
(global-set-key (kbd "C-c c") 'org-capture)
(global-set-key (kbd "C-c a") 'org-agenda)
(global-set-key (kbd "C-c b") 'org-iswitchb)

(add-hook 'org-mode-hook
          (lambda ()
            (local-set-key (kbd "C-<up>") 'org-move-subtree-up)
            (local-set-key (kbd "C-<down>") 'org-move-subtree-down)
            (local-set-key (kbd "C-c l") 'org-store-link)
            (local-set-key (kbd "C-c C-l") 'org-insert-link)))

\end{verbatim}
\subsubsection{display settings}
\label{sec:orgb65ceba}
Some config for display.
\begin{verbatim}
(setq org-hide-leading-stars t)
(setq org-log-done t)
(setq org-startup-indented t)
(setq org-startup-folded t)
(setq org-ellipsis "...")
\end{verbatim}

\subsubsection{org-habit}
\label{sec:org36a6ea3}
\begin{verbatim}
(use-package org-habit
  :delight
  :config
  (setq org-habit-graph-column 38)
  (setq org-habit-preceding-days 35)
  (setq org-habit-following-days 10)
  (setq org-habit-show-habits-only-for-today nil))
\end{verbatim}
\subsubsection{bash command}
\label{sec:org6d3b9d6}
\begin{verbatim}
(setq org-babel-sh-command "bash")
\end{verbatim}
\subsubsection{org-clock properties}
\label{sec:org2935d0e}
clock stuff into a drawer.
\begin{verbatim}
(setq org-clock-into-drawer t)
(setq org-log-into-drawer t)
(setq org-clock-int-drawer "CLOCK")
\end{verbatim}
\subsubsection{open first agenda file}
\label{sec:orga9d9959}
F12 open the first agenda file
\begin{verbatim}
(defun org-get-first-agenda-file ()
  (interactive)
  (find-file (elt org-agenda-files 0)))
(global-set-key [f12] 'org-get-first-agenda-file)
; F12 on Mac OSX displays the dashboard....
(global-set-key [C-f12] 'org-get-first-agenda-file)
\end{verbatim}
\subsubsection{org-ehtml [localhost:55555]}
\label{sec:orgd374920}
This will start serving the org files through the emacs-based webbrowser when pressing \texttt{M-f12} (on localhost:55555)
\begin{verbatim}
(use-package org-ehtml
  :ensure t
  :config
  (setq org-ehtml-docroot (expand-file-name org-directory))
  (setq org-ehtml-everything-editable t)
  (setq org-ehtml-allow-agenda t))

(defun paf/start-web-server ()
  (interactive)
  (ws-start org-ehtml-handler 55555))
(global-set-key (kbd "<M-f12>") 'paf/start-web-server)
\end{verbatim}
\subsubsection{org-link-abbrev}
\label{sec:org1076e09}
This lets one write links as e.g. [ [b:123457] ]
\begin{verbatim}
(setq org-link-abbrev-alist
      '(("b" . "http://b/")
        ("go" . "http://go/")
        ("cl" . "http://cr/")))
\end{verbatim}
\subsubsection{org-secretary}
\label{sec:org31d6123}
This package is good, but it does not do it simply. I re-modeled it
somewhat below.
\begin{verbatim}
(use-package org-secretary
  :ensure org-plus-contrib
  :config
  (setq org-sec-me "paf")
  (setq org-tag-alist '(("PRJ" . ?p)
                        ("DESIGNDOC" . ?D)
                        ("Milestone" . ?m)
                        ("DESK" . ?d)
                        ("HOME" . ?h)
                        ("VC" . ?v))))
\end{verbatim}

This is my version of the org-secretary
\begin{verbatim}
(use-package paf-secretary
  :load-path "~/Emacs"
  :bind (("\C-cw" . paf-sec-set-with)
         ("\C-cW" . paf-sec-set-where)
         ("\C-cj" . paf-sec-tag-entry))
  :config
  (setq paf-sec-me "paf")
  (setq org-tag-alist '(("PRJ" . ?p)
                        ("DESIGNDOC" . ?D)
                        ("Milestone" . ?m)
                        ("DESK" . ?d)
                        ("HOME" . ?h)
                        ("VC" . ?v))))
\end{verbatim}
\subsubsection{task tracking}
\label{sec:org37dbb7e}
Track task dependencies, and dim them in the agenda.

\begin{verbatim}
(setq org-enforce-todo-dependencies t)
(setq org-agenda-dim-blocked-tasks 'invisible)
\end{verbatim}
\subsubsection{effort \& columns mode}
\label{sec:org68fc38c}

\begin{verbatim}
(setq org-global-properties
      '(("Effort_ALL". "0 0:10 0:30 1:00 2:00 4:00 8:00 16:00")))
(setq org-columns-default-format
      "%TODO %30ITEM %3PRIORITY %6Effort{:} %10DEADLINE")
\end{verbatim}
\subsubsection{org-todo keywords}
\label{sec:orgd32b510}

\begin{verbatim}
(setq org-todo-keywords
      '((sequence "TODO(t!)" "NEXT(n!)" "STARTED(s!)" "WORKING(w!)" "|" "DONE(d!)" "CANCELLED(C@)" "DEFERRED(D@)" "SOMEDAY(S!)" "FAILED(F!)" "REFILED(R!)")
        (sequence "APPLIED(A!)" "WAITING(w!)" "ACCEPTED" "|" "REJECTED" "PUBLISHED")
        (sequence "TASK(m!)" "ACTIVE" "|" "DONE(d!)" "CANCELLED(C@)" )))

(setq org-tags-exclude-from-inheritance '("PRJ" "REGULAR")
      org-use-property-inheritance '("PRIORITY")
      org-stuck-projects '("+PRJ/-DONE-CANCELLED"
                           ;; it is considered stuck if there is no next action
                           (;"TODO"
                            "NEXT" "STARTED" "TASK") ()))

(setq org-todo-keyword-faces
      '(
        ("TODO" . (:foreground "purple" :weight bold))
        ("TASK" . (:foreground "steelblue" :weight bold))
        ("NEXT" . (:foreground "red" :weight bold))
        ("STARTED" . (:foreground "green" :weight bold))
        ("WAITING" . (:foreground "GhostWhite" :weight bold))
        ("FLAG_GATED" . (:foreground "orange" :weight bold))
        ("SOMEDAY" . (:foreground "steelblue" :weight bold))
        ("MAYBE" . (:foreground "steelblue" :weight bold))
        ("WORKING" . (:foreground "GhostWhite" :weight bold))
        ("NEW" . (:foreground "orange" :weight bold))
        ("RUNNING" . (:foreground "orange" :weight bold))
        ("WORKED" . (:foreground "green" :weight bold))
        ("FAILED" . (:foreground "red" :weight bold))
        ("REFILED" . (:foreground "gray"))
        ;; For publications
        ("APPLIED" . (:foreground "orange" :weight bold))
        ("ACCEPTED" . (:foreground "orange" :weight bold))
        ("REJECTED" . (:foreground "red" :weight bold))
        ("PUBLISHED" . (:foreground "green" :weight bold))
        ;; Other stuff
        ("ACTIVE" . (:foreground "GhostWhite" :weight bold))
        ))
\end{verbatim}
\subsubsection{org-agenda}
\label{sec:org4cd6a60}
\begin{enumerate}
\item where are my agenda files
\label{sec:org9e91408}
\begin{verbatim}
(setq org-agenda-files (list "~/org/tasks.org"))
\end{verbatim}
\item views
\label{sec:org1967422}
\begin{verbatim}
(setq org-agenda-custom-commands
      '(("t" "Hot Today" ((agenda "" ((org-agenda-span 'day)))
                          (tags-todo "-with={.+}/WAITING")
                          (tags-todo "-with={.+}+TODO=\"STARTED\"")
                          (tags-todo "/NEXT")))
        ("T" "Team Today" ((agenda "" ((org-agenda-span 'day)))
                           (tags-todo "with={.+}"
                                    ((org-super-agenda-groups
                                      '((:auto-property "with"))))
                                    )))
        ("r" "Recurring" ((tags "REGULAR")
                          (tags-todo "/WAITING")
                          (tags-todo "TODO=\"STARTED\"")
                          (tags-todo "/NEXT")))
        ("n" "Agenda and all TODO's" ((agenda "")
                                      (alltodo "")))
        ("N" "Next actions" tags-todo "-dowith={.+}/!-TASK-TODO"
         ((org-agenda-todo-ignore-scheduled t)))
        ("h" "Work todos" tags-todo "-dowith={.+}/!-TASK"
         ((org-agenda-todo-ignore-scheduled t)))
        ("H" "All work todos" tags-todo "-personal/!-TASK-CANCELLED"
         ((org-agenda-todo-ignore-scheduled nil)))
        ("A" "Work todos with doat or dowith" tags-todo
         "dowith={.+}/!-TASK"
         ((org-agenda-todo-ignore-scheduled nil)))

        ("p" "Tasks with current WITH and WHERE"
         ((tags-todo (paf-sec-replace-with-where "with={$WITH}" ".+")
                     ((org-agenda-overriding-header
                       (paf-sec-replace-with-where "Tasks with $WITH in $WHERE" "anyone" "any place"))
                      (org-super-agenda-groups
                       '((:name "" :pred paf-sec-limit-to-with-where)
                         (:discard (:anything t)))))
                     )))
        ("j" "TODO dowith and TASK with"
         ((org-sec-with-view "TODO dowith")
          (org-sec-stuck-with-view "TALK with")
          (org-sec-where-view "TODO doat")
          (org-sec-assigned-with-view "TASK with")
          (org-sec-stuck-with-view "STUCK with")
          (todo "STARTED")))
        ("J" "Interactive TODO dowith and TASK with"
         ((org-sec-who-view "TODO dowith")))))

(setq org-agenda-skip-deadline-prewarning-if-scheduled 2)
\end{verbatim}
\item delight
\label{sec:org137009d}
\begin{verbatim}
(delight 'org-agenda-mode)
\end{verbatim}
\item colors and faces
\label{sec:org3a386b0}
Make the calendar day info a bit more visible and contrasted.

\begin{verbatim}
;; Faces to make the calendar more colorful.
(custom-set-faces
 '(org-agenda-current-time ((t (:inherit org-time-grid :foreground "yellow" :weight bold))))
 '(org-agenda-date ((t (:inherit org-agenda-structure :background "pale green" :foreground "black" :weight bold))))
 '(org-agenda-date-weekend ((t (:inherit org-agenda-date :background "light blue" :weight bold)))))
\end{verbatim}
\item now marker
\label{sec:org6fd3630}
A more visible current-time marker in the agenda

\begin{verbatim}
(setq org-agenda-current-time-string ">>>>>>>>>> NOW <<<<<<<<<<")
\end{verbatim}
\item auto-refresh
\label{sec:orgf493013}

\begin{verbatim}
;; will refresh it only if already visible
(run-at-time nil 180 'update-agenda-if-visible)
;;(add-hook 'org-mode-hook
;;          (lambda () (run-at-time nil 180 'kiwon/org-agenda-redo-in-other-window)))
\end{verbatim}

This would open the agenda if any org file was opened. In the end, I
don't like this feature, so it is disabled by not tangling it.
\begin{verbatim}
;; Make this happen only if we open an org file.
(add-hook 'org-mode-hook
          (lambda () (run-with-idle-timer 600 t 'jump-to-org-agenda)))
\end{verbatim}
\item auto-save org files when idle
\label{sec:org3065b63}

This will save them regularly when the idle for more than a minute.

\begin{verbatim}
(add-hook 'org-mode-hook
    (lambda () (run-with-idle-timer 600 t 'org-save-all-org-buffers)))
\end{verbatim}
\item export
\label{sec:org5bc0491}
That's the export function to update the agenda view.
\begin{verbatim}
(setq org-agenda-exporter-settings
      '((ps-number-of-columns 2)
        (ps-portrait-mode t)
        (org-agenda-add-entry-text-maxlines 5)
        (htmlize-output-type 'font)))

(defun dmg-org-update-agenda-file (&optional force)
  (interactive)
  (save-excursion
    (save-window-excursion
      (let ((file "~/org/agenda.html"))
        (org-agenda-list)
        (org-agenda-write file)))))
\end{verbatim}
\end{enumerate}
\subsubsection{org-duration}
\label{sec:org38b7b46}
\begin{verbatim}
(use-package org-duration
  :config
  (setq org-duration-units
        `(("min" . 1)
          ("h" . 60)
          ("d" . ,(* 60 8))
          ("w" . ,(* 60 8 5))
          ("m" . ,(* 60 8 5 4))
          ("y" . ,(* 60 8 5 4 10)))
        )
  (org-duration-set-regexps))
\end{verbatim}
\subsubsection{Capture \& refile}
\label{sec:org38d9dc6}
Capture and refile stuff, with some templates that I think are useful.

Very nice post on how to get capture templats from a file:
\href{https://joshrollinswrites.com/help-desk-head-desk/org-capture-in-files/}{Org-capture
in Files}.

\begin{verbatim}
(setq org-default-notes-file (org-relative-file "inbox.org"))

(setq org-capture-templates
      `(("t" "Task"
         entry (file+headline ,(org-relative-file "inbox.org") "Tasks")
         "* TODO %?\n%U\n\n%x"
         :clock-resume t)
        ;;
        ("i" "Idea"
         entry (file+headline ,(org-relative-file "inbox.org") "Ideas")
         "* SOMEDAY %?\n%U\n\n%x"
         :clock-resume t)
        ;;
        ("m" "Meeting"
         entry (file+headline ,(org-relative-file "inbox.org") "Meetings")
         "* %?  :MTG:\n%U\n%^{with}p"
         :clock-in t
         :clock-resume t)
        ;;
        ("s" "Stand-up"
         entry (file+headline ,(org-relative-file "inbox.org") "Meetings")
         "* Stand-up  :MTG:\n%U\n\n%?"
         :clock-in t
         :clock-resume t)
        ;;
        ("1" "1:1"
         entry (file+headline ,(org-relative-file "inbox.org") "Meetings")
         "* 1:1 %^{With}  :MTG:\n%U\n:PROPERTIES:\n:with: %\\1\n:END:\n\n%?"
         :clock-in t
         :clock-resume t)
        ;;
        ("p" "Talking Point"
         entry (file+headline ,(org-relative-file "refile.org") "Talking Points")
         "* %?  :TALK:\n%U\n%^{dowith}p"
         :clock-keep t)
        ;;
        ("j" "Journal"
         entry (file+olp+datetree ,(org-relative-file "journal.org"))
         "* %?\n%U"
         :clock-in t
         :clock-resume t
         :kill-buffer t)))

;; show up to 2 levels for refile targets, in all agenda files
(setq org-refile-targets '((org-agenda-files . (:maxlevel . 2))))
(setq org-log-refile t)  ;; will add timestamp when refiled.

;; from: http://doc.norang.ca/org-mode.html
;; Exclude DONE state tasks from refile targets
(defun bh/verify-refile-target ()
  "Exclude todo keywords with a done state from refile targets"
  (not (member (nth 2 (org-heading-components)) org-done-keywords)))
(setq org-refile-target-verify-function 'bh/verify-refile-target)
\end{verbatim}
\subsubsection{OrgRoam templates}
\label{sec:orgad8abf7}
\begin{verbatim}
(setq org-roam-capture-templates
      `(("m" "Meeting" entry (function org-roam--capture-get-point)
             "* %?\n%U\n%^{with}\n"
             :file-name "meeting/%<%Y%m%d%H%M%S>-${slug}"
             :head "#+title: ${title}\n#+roam_tags: %^{with}\n\n"
             )))

\end{verbatim}
\subsubsection{org-babel}
\label{sec:orgb053a25}
What kind of code block languages do I need

\begin{verbatim}
(setq org-confirm-babel-evaluate 'nil) ; Don't ask before executing

(org-babel-do-load-languages
 'org-babel-load-languages
 '(
   (R . t)
   (dot . t)
   (emacs-lisp . t)
   (gnuplot . t)
   (python . t)
   (ledger . t)
   ;;(sh . t)
   (latex . t)
   (shell . t)
  ))
\end{verbatim}
\subsubsection{org-export}
\label{sec:org9deeb59}
Add a few formats to the export functionality of org-mode.

\begin{verbatim}
(use-package ox-odt
  :defer)
(use-package ox-taskjuggler
  :defer)
(use-package ox-impress-js
  :defer)
\end{verbatim}
\subsubsection{plant-uml}
\label{sec:orgeb312a3}

Tell where PlantUML is to be found. This needs to be downloaded and
installed separately, see the \href{http://plantuml.com/}{PlantUML website}.

You could install the PlantUML JAR file with this snippet:

\begin{verbatim}
# Get a version of the PlantUML jar file.
install_pkg -x wget wget

URL='http://sourceforge.net/projects/plantuml/files/plantuml.jar/download'
DIR="${HOME}/Emacs"
if [[ ! -e "${DIR}/plantuml.jar" ]]; then
    [[ -d "${DIR}" ]] || mkdir -p "${DIR}"
    (cd "${DIR}" && wget -O plantuml.jar "${URL}")
    ls -l "${DIR}/plantuml.jar"
fi
\end{verbatim}

\begin{verbatim}
(use-package plantuml-mode
 :ensure t
 :config
  (setq plantuml-jar-path "~/workspace/me/myblog/content/downloads/plantuml.jar")
  (setq org-plantuml-jar-path "~/workspace/me/myblog/content/downloads/plantuml.jar")
  ;; Let us edit PlantUML snippets in plantuml-mode within orgmode
  (add-to-list 'org-src-lang-modes '("plantuml" . plantuml))
  ;; make it load this language (for export ?)
  (org-babel-do-load-languages 'org-babel-load-languages '((plantuml . t)))
  ;; Enable plantuml-mode for PlantUML files
  (add-to-list 'auto-mode-alist '("\\.plantuml\\'" . plantuml-mode)))
\end{verbatim}

\subsubsection{PDF-Tools}
\label{sec:org3e536de}
A bit difficult to find the docs of how to use it, but it seems quite
useful.

Disabled, as it causes only trouble to me, and I am not really using
it anyway.

\begin{verbatim}
(use-package pdf-tools
  :if (and (eq system-type 'gnu/linux)  ;; Set it up on Linux
           (not (string-prefix-p "aarch64" system-configuration)))  ;; but not mobile devices
  :pin manual  ;; update only manually
  :config
  ;; initialize
  (pdf-tools-install)
  (setq-default pdf-view-display-size 'fit-page)           ;; Fit to page when opening
  (add-hook 'pdf-view-mode-hook (lambda () (cua-mode 0)))  ;; turn off cua so copy works
  (setq pdf-view-resize-factor 1.1)                        ;; more fine-grained zoom control
  ;; keyboard shortcuts
  (define-key pdf-view-mode-map (kbd "h") 'pdf-annot-add-highlight-markup-annotation)
  (define-key pdf-view-mode-map (kbd "t") 'pdf-annot-add-text-annotation)
  (define-key pdf-view-mode-map (kbd "D") 'pdf-annot-delete))

(use-package org-pdfview
  :after (pdf-tools)
  :init
  (add-to-list 'org-file-apps '("\\.pdf\\'" . org-pdfview-open))
  (add-to-list 'org-file-apps '("\\.pdf::\\([[:digit:]]+\\)\\'" . org-pdfview-open)))
\end{verbatim}

\begin{verbatim}
# For all the native apps related to PDF tools
# I did not sintall it on Max OSX yet.
if [[ "$(uname)" != "Darwin" ]]; then
  install_pkg pdf-tools
fi
\end{verbatim}
\subsubsection{yankpad}
\label{sec:orgd228adf}
Check out the \href{https://kungsgeten.github.io/yankpad.html}{blog post} (and the \href{https://kungsgeten.github.io/yankpad13.html}{follow-up}) and the \href{https://github.com/Kungsgeten/yankpad}{package docs}.

\begin{verbatim}
(use-package yankpad
  :ensure t
  :defer
  :init
  (setq yankpad-file (org-relative-file "yankpad.org"))
  :config
  (bind-key "C-c f y m" 'yankpad-map)
  (bind-key "C-c f y e" 'yankpad-expand))
\end{verbatim}

\subsubsection{\href{https://www.eliasstorms.net/zetteldeft/}{Zetteldeft}}
\label{sec:org0cb593c}
This is a note-taking packages inspired by the principles of the
\href{https://zettelkasten.de/}{Zettelkasten}

\begin{verbatim}
(use-package deft
  :ensure t)
(use-package avy
  :ensure t)

(use-package zetteldeft
  :ensure t
  :after (org deft avy)

  :config
  (setq deft-extensions '("org" "md" "txt"))
  (setq deft-directory (org-relative-file "Zettelkasten"))
  (setq deft-recursive t)

  :bind (("C-c z d" . deft)
         ("C-c z D" . zetteldeft-deft-new-search)
         ("C-c z R" . deft-refresh)
         ("C-c z s" . zetteldeft-search-at-point)
         ("C-c z c" . zetteldeft-search-current-id)
         ("C-c z f" . zetteldeft-follow-link)
         ("C-c z F" . zetteldeft-avy-file-search-ace-window)
         ("C-c z l" . zetteldeft-avy-link-search)
         ("C-c z t" . zetteldeft-avy-tag-search)
         ("C-c z T" . zetteldeft-tag-buffer)
         ("C-c z i" . zetteldeft-find-file-id-insert)
         ("C-c z I" . zetteldeft-find-file-full-title-insert)
         ("C-c z o" . zetteldeft-find-file)
         ("C-c z n" . zetteldeft-new-file)
         ("C-c z N" . zetteldeft-new-file-and-link)
         ("C-c z r" . zetteldeft-file-rename))
)
\end{verbatim}

Update the version by downloading the latest version here:

\begin{verbatim}
wget https://raw.githubusercontent.com/EFLS/zetteldeft/master/zetteldeft.el -O ~/Emacs/zetteldeft.el
\end{verbatim}

\subsubsection{reload org}
\label{sec:org3020721}
Ran into a \href{https://github.com/seagle0128/.emacs.d/issues/129}{strange error}, and reloading org at the end is the
solution:

\begin{verbatim}
(org-reload)
\end{verbatim}

\section{Write chinese}
\label{sec:orga42552f}
I have been using this one w/ reasonable success.
\begin{verbatim}
(use-package pyim
  :ensure
  :defer 30
  :config
  ;; 激活 basedict 拼音词库
  (use-package pyim-basedict
    :ensure nil
    :config (pyim-basedict-enable))

  ;; 五笔用户使用 wbdict 词库
  ;; (use-package pyim-wbdict
  ;;   :ensure nil
  ;;   :config (pyim-wbdict-gbk-enable))

  (setq default-input-method "pyim")

  ;; 我使用全拼
  (setq pyim-default-scheme 'quanpin)

  ;; 设置 pyim 探针设置，这是 pyim 高级功能设置，可以实现 *无痛* 中英文切换 :-)
  ;; 我自己使用的中英文动态切换规则是：
  ;; 1. 光标只有在注释里面时，才可以输入中文。
  ;; 2. 光标前是汉字字符时，才能输入中文。
  ;; 3. 使用 M-j 快捷键，强制将光标前的拼音字符串转换为中文。
  (setq-default pyim-english-input-switch-functions
                '(pyim-probe-dynamic-english
                  pyim-probe-isearch-mode
                  pyim-probe-program-mode
                  pyim-probe-org-structure-template))

  (setq-default pyim-punctuation-half-width-functions
                '(pyim-probe-punctuation-line-beginning
                  pyim-probe-punctuation-after-punctuation))

  ;; 开启拼音搜索功能
  ;; (pyim-isearch-mode 1)

  ;; 使用 pupup-el 来绘制选词框
  (setq pyim-page-tooltip 'popup)

  ;; 选词框显示5个候选词
  (setq pyim-page-length 5)

  ;; 让 Emacs 启动时自动加载 pyim 词库
  (add-hook 'emacs-startup-hook
            #'(lambda () (pyim-restart-1 t)))
  :bind
  (;与 pyim-probe-dynamic-english 配合
  ("M-j" . pyim-convert-code-at-point)

   ("C-;" . pyim-delete-word-from-personal-buffer)))
   (setq default-input-method "pyim")
   (global-set-key (kbd "C-\\") 'toggle-input-method)
\end{verbatim}
\section{Office stuff}
\label{sec:org2cd8ff1}
Using emacs in a corporate env can be daunting.

\subsection{mu4e email}
\label{sec:orgff3fcd7}
mu4e only a emac client. The workhorse are davmail, isync and mu. You need to install these offline and hook them up using the config below.

\subsubsection{install \& minimal setup}
\label{sec:org55aa375}
First thing first, load the package. As you can see, I have used \texttt{apt-get} install mu4e.

\begin{verbatim}
(add-to-list 'load-path "/usr/share/emacs/site-lisp/mu4e/")
(require 'mu4e)
\end{verbatim}

Then we load \texttt{maildirs-extension}:

\begin{verbatim}
(use-package mu4e-maildirs-extension
  :ensure
  :config)
(mu4e-maildirs-extension)
\end{verbatim}

Now, we tell emacs I want to use \texttt{mu4e} as email client:

\begin{verbatim}
(setq mail-user-agent 'mu4e-user-agent)
\end{verbatim}

\subsubsection{how to get mails}
\label{sec:org221baae}
Setup location of my maildir.
\begin{verbatim}
(setq mu4e-maildir (expand-file-name "~/Maildir"))
\end{verbatim}

Sync email by calling \texttt{mbsync}:
\begin{verbatim}
(setq mu4e-get-mail-command "mbsync -a")
\end{verbatim}

How often should I check? Value in seconds:
\begin{verbatim}
(setq mu4e-update-interval 600)
\end{verbatim}

Setup some common folders:
\begin{verbatim}
(setq mu4e-drafts-folder "/drafts"
      mu4e-sent-folder   "/sent"
      mu4e-trash-folder  "/trash")
\end{verbatim}

Setup some shortcuts as bookmarks:
\begin{verbatim}
(add-to-list
 'mu4e-bookmarks
 '("flag:attach"
   "Messages with attachment"
   ?a) t)

(add-to-list
 'mu4e-bookmarks
 '("size:5M..500M"
   "Big messages"
   ?b) t)

(add-to-list
 'mu4e-bookmarks
 '("flag:flagged"
   "Flagged messages"
   ?f) t)
\end{verbatim}

\subsubsection{list view}
\label{sec:org23b6502}
Customize the list view header:
\begin{verbatim}
(setq mu4e-headers-date-format "%b-%d %a"
      mu4e-headers-fields '((:date . 10)
                            (:flags . 5)
                            (:recipnum . 3)
                            (:from-or-to . 10)
                            (:thread-subject . nil)))
\end{verbatim}

Skip duplicates:

\begin{verbatim}
(setq mu4e-headers-skip-duplicates t)
\end{verbatim}

Showing related in a tree fashion so I know the context:

\begin{verbatim}
(setq mu4e-headers-include-related t)
\end{verbatim}

Here is a fun one. I noticed that the email thread grows like that greedy snake game, a pretty good sign that the team is malfunctioning :) So we add a displayed number on the number of recipients, and just watch it grow:
\begin{verbatim}
(add-to-list 'mu4e-header-info-custom
             '(:recipnum .
                         ( :name "Number of recipients"  ;; long name, as seen in the message-view
                                 :shortname "R#"           ;; short name, as seen in the headers view
                                 :help "Number of recipients for this message" ;; tooltip
                                 :function (lambda (msg)
                                             (format "%d"
                                                     (+ (length (mu4e-message-field msg :to))
                                                        (length (mu4e-message-field msg :cc))))))))

\end{verbatim}
\subsubsection{read \& attachment}
\label{sec:org626d795}
Save attachment to:

\begin{verbatim}
(setq mu4e-attachment-dir "~/Downloads")
\end{verbatim}

To open HTML in browser:
\begin{verbatim}
(add-to-list 'mu4e-view-actions
             '("ViewInBrowser" . mu4e-action-view-in-browser)
             t)
\end{verbatim}

Attempt to show images when viewing messages:

\begin{verbatim}
(setq mu4e-view-show-images t)
\end{verbatim}

Use imagemagick, if available:

\begin{verbatim}
(when (fboundp 'imagemagick-register-types)
  (imagemagick-register-types))
\end{verbatim}
\subsubsection{write new}
\label{sec:org4183f51}
My signature line:

\begin{verbatim}
(setq mu4e-compose-signature
      (concat
       "Best regards,\n\n"
       "Feng Xia\n"
       "W: http://www.lenovo.com\n"))
\end{verbatim}
\subsubsection{send}
\label{sec:orge3a7cac}
Sending email is to use SMTP. Here I'm showing to use the company one
through \texttt{davmail} which runs smtp listen on port \texttt{1025}. You can also
use Gmail/hotmail:

\begin{verbatim}
(use-package smtpmail
  :ensure t
  :config
  (setq send-mail-function 'smtpmail-send-it
        user-mail-address "fxia1@lenovo.com"
        smtpmail-debug-info t
        smtpmail-smtp-user "fxia1@lenovo.com"
        smtpmail-default-smtp-server "localhost"
        smtpmail-auth-credentials (expand-file-name "~/.authinfo")
        smtpmail-smtp-service 1025
        smtpmail-stream-type nil
        starttls-use-gnutls nil
        starttls-extra-arguments nil))
\end{verbatim}

There is a toggle where you can queue the sending email first instead
of sending to server immediately. This is useful for travel when you
just write, then queue, then when reconnected to the internect, send.

\begin{verbatim}
(setq smtpmail-queue-mail nil
      smtpmail-queue-dir "~/Maildir/queue/cur")
\end{verbatim}
\subsubsection{reply}
\label{sec:org04a3b07}
When replying an email, auto fill in my info:
\begin{verbatim}
(setq mu4e-compose-reply-to-address "fxia1@lenovo.com"
      user-mail-address "fxia1@lenovo.com"
      user-full-name "Feng Xia"
      message-signature  (concat
                          "Feng Xia\n"
                          "http://snapshots.world.s3-website.ap-northeast-2.amazonaws.com/\n")
      message-citation-line-format "On %Y-%m-%d %H:%M:%S, %f wrote:"
      message-citation-line-function 'message-insert-formatted-citation-line
      mu4e-headers-results-limit 500)
\end{verbatim}
\subsubsection{kill all buffer upon exit}
\label{sec:org09b35f4}
mu4e can open a lot of writing buffers. Just kill them all when we exit mu4e:

\begin{verbatim}
(setq message-kill-buffer-on-exit t)
\end{verbatim}
\subsubsection{write html}
\label{sec:org844e116}
Emacs doesn't like html email body. Nor do I. But one thing I found
out is that markdown mode table will look terrible as plain text on
receiving end. So we now write email in org mode, then use this to
convert it to html before sending:
\begin{verbatim}
(use-package org-mime
  :ensure
  :config
  (setq org-mime-export-ascii 'utf-8))
\end{verbatim}

How about writing it in markdown and export?
\begin{verbatim}
(defun multipart-html-message (plain html)
  "Creates a multipart HTML email with a text part and an html part."
  (concat "<#multipart type=alternative>\n"
          "<#part type=text/plain>"
          plain
          "<#part type=text/html>\n"
          html
          "<#/multipart>\n"))

(defun convert-message-to-markdown ()
  "Convert the message in the current buffer to a multipart HTML email.

The HTML is rendered by treating the message content as Markdown."
  (interactive)
  (unless (executable-find "pandoc")
    (error "Pandoc not found, unable to convert message"))
  (let* ((begin
          (save-excursion
            (goto-char (point-min))
            (search-forward mail-header-separator)))
         (end (point-max))
         (html-buf (generate-new-buffer "*mail-md-output*"))
         (exit-code
          (call-process-region begin end "pandoc" nil html-buf nil
                               "--quiet" "-f" "gfm" "-t" "html"))
         (html (format "<html>\n<head></head>\n<body>\n%s\n</body></html>\n"
                (with-current-buffer html-buf
                  (buffer-substring (point-min) (point-max)))))
         (raw-body (buffer-substring begin end)))
    (when (not (= exit-code 0))
      (error "Markdown conversion failed, see %s" (buffer-name html-buf)))
    (with-current-buffer html-buf
      (set-buffer-modified-p nil)
      (kill-buffer))
    (undo-boundary)
    (delete-region begin end)
    (save-excursion
      (goto-char begin)
      (newline)
      (insert (multipart-html-message raw-body html)))))

(defun message-md-send (&optional arg)
  "Convert the current buffer and send it.
If given prefix arg ARG, skips markdown conversion."
  (interactive "P")
  (unless arg
    (convert-message-to-markdown))
  (message-send))

(defun message-md-send-and-exit (&optional arg)
  "Convert the current buffer and send it, then exit from mail buffer.
If given prefix arg ARG, skips markdown conversion."
  (interactive "P")
  (unless arg
    (convert-message-to-markdown))
  (message-send-and-exit))

(with-eval-after-load 'message
 (define-key message-mode-map (kbd "C-c C-s") #'message-md-send)
 (define-key message-mode-map (kbd "C-c C-c") #'message-md-send-and-exit))
\end{verbatim}
\section{Hobby stuff}
\label{sec:org19a3505}
Things I do as hobby and could now do them inside emacs.

\section{Learn org}
\label{sec:org8416cbf}
Here is my sandbox that I'm gonna learn writing org to manage my tasks and life.
\subsection{write todo}
\label{sec:orgb117d75}
\subsubsection{{\bfseries\sffamily WORKING} create a new task?}
\label{sec:org066dc7e}
Add something content to this todo, and you can write as much as you'd
like so to describe your idea well with \textbf{bold} and \emph{slant} and \uline{underscore} and \texttt{verbatim}?

\begin{center}
\begin{tabular}{ll}
this & that\\
\hline
how about a table & like this?\\
\end{tabular}
\end{center}

Then I can write more and more and more.
\subsubsection{{\bfseries\sffamily TODO} The, move on to the next one \textasciitilde{}\textasciitilde{} this is really neat!}
\label{sec:orgff5d0a3}
\subsubsection{{\bfseries\sffamily WAITING} some task}
\label{sec:org93930bc}
\subsubsection{{\bfseries\sffamily TODO} try another todo}
\label{sec:orgc098295}
\subsubsection{{\bfseries\sffamily TODO} lab cleanup}
\label{sec:orge9c37a8}
\begin{enumerate}
\item switch login:
\begin{itemize}
\item (admin, SYSadmin2019)
\item (admin,admin)
\item (admin, CPadmin2019)
\item (admin, LOCadmin2020)
\end{itemize}
\item switch serial connection: \texttt{do screen /dev/ttyUSB0 115200,cs8}. For older switches, baud rate is 9600.
\end{enumerate}


\begin{center}
\begin{tabular}{rlrll}
 & Device & IP & pwd & In excel?\\
\hline
1 & R1U40 & 10.240.43.51 &  & \\
2 & R1U39 & 10.240.41.146 &  & \\
3 & R1U38 &  &  & \\
4 & R1U37 & 10.240.41.147 & LOCadmin2020 & \\
5 & R31U41 & 192.168.50.50 & CPadmin2019 & n/a\\
6 & R20U40 & 10.240.20.140 & CPadmin2019 & no\\
7 & R0U40 & 10.240.40.84 &  & \\
8 & ROU39 & 10.240.42.204 &  & \\
9 & ROU38 & 10.240.42.205 &  & \\
10 & ROU37 & 10.240.42.158 &  & \\
11 & ROU36 & 10.240.40.83 &  & \\
12 & ROU35 & 10.240.40.195 &  & no\\
13 & R2U39 & 10.240.41.144 &  & \\
14 & R2U38 & 10.240.43.54 & admin & \\
15 & R2U37 & 10.240.41.145 &  & \\
16 & R3U40 & 10.240.41.99 &  & \\
17 & R3U39 & 10.240.41.100 &  & \\
18 & R3U38 & 10.240.41.101 &  & \\
19 & R4U39 & 10.240.43.28 &  & \\
20 & R4U37 & 10.240.43.29 &  & \\
21 & R5U39 & 10.240.41.103 &  & no\\
22 & R5U38 & 10.240.42.228 &  & \\
23 & R5U37 & 10.240.124.10 &  & \\
24 & R5U36 & 10.240.40.248 &  & \\
25 & R19U40 & 10.240.20.147 & CPadmin2019 & \\
26 & R19U25 & 10.240.20.249 & CPadmin2019 & \\
27 & R21U41 & 10.240.20.92 & CPadmin2019 & \\
28 & R21U40 & 10.240.20.93 & CPadmin2019 & \\
29 & R22U40 & 10.240.18.30 & CPadmin2019 & \\
30 & R22U39 & 10.240.20.49 & CPadmin2019 & \\
31 & R22U38 & 10.240.22.10 & CPadmin2019 & \\
32 & R23U40 & 10.240.20.94 & CPadmin2019 & \\
33 & R23U39 &  &  & \\
34 & R23U38 &  &  & \\
35 & R23U37 & 10.240.20.103 & CPadmin2019 & \\
36 & R23U36 & 10.240.18.12 & CPadmin2019 & \\
37 & R23U34 & 10.240.20.97 & none & no\\
38 & R24U41 & 10.240.20.141 & CPadmin2019 & \\
39 & R24U40 & 10.240.20.138 & CPadmin2019 & \\
40 & R25U41 & 10.240.20.142 & CPadmin2019 & \\
41 & R25U40 & 10.240.20.139 & CPadmin2019 & \\
42 & R25U35 & 10.240.20.109 & none & no\\
\end{tabular}
\end{center}

\begin{enumerate}
\item Active BMCs
\begin{itemize}
\item 10.240.41.71
\item 10.240.41.74
\item 10.240.20.148
\item 10.240.41.248
\item 10.240.41.249
\item 10.240.41.250
\item 10.240.41.251
\item 10.240.41.246
\item 10.240.41.247
\item 10.240.41.243
\item 10.240.41.229
\item 10.240.41.228
\item 10.240.43.49
\end{itemize}
\item
\end{enumerate}
\end{document}
